% ===================================================================
%  EN47  Landweber《K 理论与椭圆算子》中文译注本  汇总文件
%  Lead 独占文件：译者不要修改本文件
% ===================================================================
% ===================================================================
%  EN47  Landweber《K 理论与椭圆算子》中文译注本  导言区
%  Lead 独占文件：译者不要修改本文件
% ===================================================================
\documentclass[11pt,a4paper]{article}

\usepackage[UTF8]{ctex}
\usepackage{amsmath,amssymb,amsthm,mathtools}
\usepackage{mathrsfs}
\usepackage{bm}
\usepackage{graphicx}
\usepackage{tikz}
\usepackage{tikz-cd}
\usepackage[a4paper,top=2.4cm,bottom=2.6cm,left=2.4cm,right=2.4cm]{geometry}
\usepackage[dvipsnames]{xcolor}
\usepackage{enumitem}
\usepackage{array}
\usepackage{multicol}
\usepackage{fancyhdr}
\usepackage{caption}
\usepackage{titlesec}
\usepackage[hidelinks,bookmarksnumbered]{hyperref}

% ---------- 版面 ----------
\linespread{1.18}
\setlength{\parindent}{2em}
\setlength{\parskip}{0.25em}
\captionsetup{font=small}

% ---------- 页眉页脚 ----------
\pagestyle{fancy}
\fancyhf{}
\fancyhead[L]{\small K 理论与椭圆算子}
\fancyhead[R]{\small 中文译注本}
\fancyfoot[C]{\small\thepage}
\renewcommand{\headrulewidth}{0.3pt}

% ---------- 节标题 ----------
\newcommand{\seclabel}{第 \thesection{} 节}
\titleformat{\section}{\Large\bfseries}{\seclabel}{1em}{}
\titleformat{\subsection}{\large\bfseries}{\thesubsection}{0.8em}{}
\titlespacing*{\section}{0pt}{2.2ex plus 1ex minus .2ex}{1.2ex plus .2ex}
\titlespacing*{\subsection}{0pt}{1.8ex plus 1ex minus .2ex}{0.8ex plus .2ex}

% ---------- 定理类环境（原文全部不编号） ----------
\theoremstyle{definition}
\newtheorem*{definition}{定义}
\newtheorem*{theorem}{定理}
\newtheorem*{lemma}{引理}
\newtheorem*{proposition}{命题}
\newtheorem*{claim}{断言}
\newtheorem*{example}{例}
\newtheorem*{corollary}{推论}
\newtheorem*{problem}{问题}
\newtheorem*{remark}{注}
\renewcommand{\proofname}{证明}
\renewcommand{\qedsymbol}{$\square$}
\renewcommand{\refname}{参考文献}

% ---------- 译者注 ----------
\newcommand{\transnote}[1]{%
  \par\medskip\noindent
  {\small\color{Sepia}\textbf{【译者注】}#1\par}\medskip}

% ---------- 常用记号 ----------
\newcommand{\R}{\mathbb{R}}
\newcommand{\C}{\mathbb{C}}
\newcommand{\Z}{\mathbb{Z}}
\newcommand{\Q}{\mathbb{Q}}
\newcommand{\N}{\mathbb{N}}
\newcommand{\PP}{\mathbb{P}}
\newcommand{\KK}{\mathbb{K}}
\newcommand{\OO}{\mathcal{O}}
\newcommand{\GG}{\mathcal{G}}
\newcommand{\HH}{\mathcal{H}}
\newcommand{\MM}{\mathcal{M}}
\newcommand{\LL}{\mathcal{L}}
\newcommand{\de}{\mathrm{d}}
\newcommand{\dd}{\mathrm{d}}
\newcommand{\id}{\mathrm{id}}
\newcommand{\pt}{\mathrm{pt}}
\newcommand{\rk}{\operatorname{rk}}
\newcommand{\tr}{\operatorname{tr}}
\newcommand{\sgn}{\operatorname{sgn}}
\newcommand{\vol}{\operatorname{vol}}
\newcommand{\supp}{\operatorname{supp}}
\newcommand{\Vect}{\operatorname{Vect}}
\newcommand{\Coker}{\operatorname{Coker}}
\newcommand{\Ker}{\operatorname{Ker}}
\newcommand{\Ind}{\operatorname{Index}}
\newcommand{\tind}{\operatorname{t\text{-}ind}}
\newcommand{\aind}{\operatorname{a\text{-}ind}}
\newcommand{\Td}{\operatorname{Td}}
\newcommand{\ch}{\operatorname{ch}}
\newcommand{\Hom}{\operatorname{Hom}}
\newcommand{\End}{\operatorname{End}}
\newcommand{\Aut}{\operatorname{Aut}}
\newcommand{\GL}{\operatorname{GL}}
\newcommand{\U}{\mathrm{U}}
\newcommand{\SU}{\mathrm{SU}}
\newcommand{\SO}{\mathrm{SO}}
\newcommand{\Sp}{\mathrm{Sp}}
\newcommand{\Spin}{\mathrm{Spin}}
\newcommand{\IndexClass}{\Im}

% 花体字母
\newcommand{\cA}{\mathcal{A}}\newcommand{\cB}{\mathcal{B}}
\newcommand{\cC}{\mathcal{C}}\newcommand{\cD}{\mathcal{D}}
\newcommand{\cE}{\mathcal{E}}\newcommand{\cF}{\mathcal{F}}
\newcommand{\cG}{\mathcal{G}}\newcommand{\cH}{\mathcal{H}}
\newcommand{\cI}{\mathcal{I}}\newcommand{\cJ}{\mathcal{J}}
\newcommand{\cK}{\mathcal{K}}\newcommand{\cL}{\mathcal{L}}
\newcommand{\cM}{\mathcal{M}}\newcommand{\cN}{\mathcal{N}}
\newcommand{\cO}{\mathcal{O}}\newcommand{\cP}{\mathcal{P}}
\newcommand{\cQ}{\mathcal{Q}}\newcommand{\cR}{\mathcal{R}}
\newcommand{\cS}{\mathcal{S}}\newcommand{\cT}{\mathcal{T}}
\newcommand{\cU}{\mathcal{U}}\newcommand{\cV}{\mathcal{V}}
\newcommand{\cW}{\mathcal{W}}\newcommand{\cX}{\mathcal{X}}
\newcommand{\cY}{\mathcal{Y}}\newcommand{\cZ}{\mathcal{Z}}
% 哥特体字母
\newcommand{\fA}{\mathfrak{A}}\newcommand{\fB}{\mathfrak{B}}
\newcommand{\fC}{\mathfrak{C}}\newcommand{\fD}{\mathfrak{D}}
\newcommand{\fE}{\mathfrak{E}}\newcommand{\fF}{\mathfrak{F}}
\newcommand{\fG}{\mathfrak{G}}\newcommand{\fH}{\mathfrak{H}}
\newcommand{\fI}{\mathfrak{I}}\newcommand{\fJ}{\mathfrak{J}}
\newcommand{\fK}{\mathfrak{K}}\newcommand{\fL}{\mathfrak{L}}
\newcommand{\fM}{\mathfrak{M}}\newcommand{\fN}{\mathfrak{N}}
\newcommand{\fO}{\mathfrak{O}}\newcommand{\fP}{\mathfrak{P}}
\newcommand{\fQ}{\mathfrak{Q}}\newcommand{\fR}{\mathfrak{R}}
\newcommand{\fS}{\mathfrak{S}}\newcommand{\fT}{\mathfrak{T}}
\newcommand{\fU}{\mathfrak{U}}\newcommand{\fV}{\mathfrak{V}}
\newcommand{\fW}{\mathfrak{W}}\newcommand{\fX}{\mathfrak{X}}
\newcommand{\fY}{\mathfrak{Y}}\newcommand{\fZ}{\mathfrak{Z}}

% ---------- 杂项 ----------
\newcommand{\eps}{\varepsilon}
\newcommand{\BT}{\widehat{\otimes}}   % 原文的外积符号实为 \widehat{\otimes}（非 \boxtimes）
\newcommand{\iso}{\cong}
\newcommand{\into}{\hookrightarrow}
\newcommand{\onto}{\twoheadrightarrow}
\newcommand{\restr}[2]{{#1}\big|_{#2}}
\newcommand{\ip}[2]{\langle #1,#2\rangle}
\allowdisplaybreaks[2]
\emergencystretch=2em


\begin{document}

\begin{center}
  {\LARGE\bfseries K 理论与椭圆算子}\\[0.6em]
  {\large 中文译注本}\\[1.2em]
  {\large Gregory D.~Landweber}\\[0.3em]
  {\small University of Oregon, Mathematics Department}\\[1.2em]
  {\small 底本：\emph{K-Theory and Elliptic Operators}, arXiv:math/0504555v1 [math.AT], 27 April 2005}
\end{center}

\vspace{0.6em}
\noindent\rule{\textwidth}{0.4pt}
\begin{minipage}{\textwidth}
\small
\textbf{译本说明}\quad
本译本的底本为 Gregory D.~Landweber 的讲义 \emph{K-Theory and Elliptic Operators}
（arXiv:math/0504555v1）。原书 50 页，其中第 1 页为 arXiv 扉页、第 2 页为目录，
正文自第 3 页起；本译本的页码与原书页码一一对应，正文中每页起首以边注形式的
注释标出\texttt{\% ===== 原书 p.N =====}，便于对照。原书全部公式与定理均无编号，
译本亦不编号；节号、节标题的语义与环境类型保持不变。译文中的
\textbf{【译者注】}为译者所加，只分三类：易错预警、含糊点拨、背景直觉，
不代替读者阅读原文，也不替作者补证明。人名与标准术语保留英文原名，
术语首次出现时给出英文原词。
\end{minipage}
\noindent\rule{\textwidth}{0.4pt}

\vspace{1.2em}
\noindent\textbf{内容提要}\quad
本讲义是拓扑 K 理论与 Atiyah–Singer 指标定理的入门读本，对象是已学过代数拓扑的
研究生或高年级本科生。书中大部分材料取自 Atiyah 的经典著作《K-Theory》以及他与
Singer 合作的系列论文《The Index of Elliptic Operators》。此外还包括等变 K 理论、
G-指标定理，以及 Bott 的论文《The Index Theorem for Homogeneous Differential
Operators》中的内容。附录 A 给出 Bott 周期性的一个证明，并利用椭圆算子族的指标，
概述 K 理论 Thom 同构定理的通常形式与等变形式的证明；附录 B 不使用同调代数知识，
直接导出 Atiyah–Hirzebruch 谱序列。本书最初是作者在普林斯顿读本科时一系列讲座的
笔记，现在这个形式则用于俄勒冈大学的研究生课程。

\vspace{0.8em}
\noindent{\small 2000 年数学主题分类：主分类 55N15, 58J20；次分类 19L47}

\vspace{0.4em}
\noindent{\small 关键词：拓扑 K 理论，Atiyah–Singer 指标定理，Bott 周期性}

\vspace{1.4em}
\tableofcontents
\newpage

\setcounter{section}{-1}
% ===== 原书 p.3 =====
\section{指标问题}

\subsection{引言}

本文的中心议题是 Atiyah--Singer 指标定理。这是二十世纪后半叶最重要的数学成果之一，它在代数拓扑、微分几何与分析之间建立起根本的联系。具体地说，指标定理把线性椭圆微分算子的指标用若干拓扑不变量表示出来。寻找这种表示的问题最初由 Israel Gel'fand 于 1960 年提出，1962 年秋，Michael Atiyah 与 Isadore Singer 在哈佛与 MIT 解决了它。1963 年初，Atiyah 与 Singer 在《Bulletin of the American Mathematical Society》上联名发表论文 [9]，首次宣布指标定理，其中用配边论把指标表示成示性类的形式。1968 年，Atiyah 与 Singer 在《Annals of Mathematics》上以三篇论文的形式发表了指标定理的第二个证明（该系列的第二篇实际出自 Atiyah 与 Graeme Segal 之手）。第二个证明用 K 理论重新表述了指标定理，完全避开了上同调与配边论，并把它推广为等变的 $G$-指标定理与椭圆算子族的版本。1968 年以来，指标定理出现了多种不同的证明（例如热方程证明与超对称证明），它在数学物理中也变得重要起来。

在本文中，我们从第 1 节开始讨论 K 理论，这是指标定理的基本语言。K 理论主要由 Atiyah 与 Hirzebruch 在 1960 年代初发展起来；一般性的参考文献是 Atiyah 的经典著作 [6]。第 1 节末尾简要回顾示性类，以及 K 理论与上同调的关系。第 2、3 节沿 [10] 与 [11] 的思路给出指标定理的证明，其中第 2 节稍作停顿，介绍指标定理的上同调形式并给出几个简单应用。计入紧 Lie 群 $G$ 的作用之后的等变推广放在第 4 节。等变 K 理论由 Atiyah 与 Segal 在牛津发展起来，完整的处理见 Segal 的 [17]。在第 4 节中，我们先陈述 $G$-指标定理，然后按 Raoul Bott 的论文 [12] 考察齐性微分算子这一特殊情形。第一个附录利用一族微分算子的指标给出 Bott 周期性定理的一个证明；这部分讨论基本沿 Atiyah 的论文 [5]，并把该技巧推广，用来证明 Thom 同构定理的通常形式与等变形式。第二个附录不假定任何同调代数知识，直接导出 Atiyah--Hirzebruch 谱序列。

书中的材料并非原创；事实上，所有结果在作者出生之前就已为人所知。本书要做的是把 1963 至 1968 年这一关键时期 K 理论与指标理论各方面的发展综合起来，重点放在 Atiyah 及其各位合作者的工作上。书中试图统一这一领域各篇论文的记号，并详细算出若干例子。特别是，第 2 节对奇维流形上指标的处理，以及附录中关于粘合函数的例子，都比现有文献详细得多。

本书是为高年级本科生与研究生写的指标定理导论，假定他们几乎没有或完全没有 K 理论与指标理论的基础。不过，书中假定读者熟悉向量丛的基本理论；若不熟悉，应先浏览 [16] 的前几节或 [6] 的第 1 节，再读本书。建议读者具备较好的上同调论基础并了解一些示性类，但对指标定理本身的证明而言，这些并非严格必需。第 4 节还需要一点表示论的知识，可参阅 [1]。

% ===== 原书 p.4 =====
\subsection{定义与例子}

设 $X$ 是紧光滑流形。给定 $X$ 上的光滑复向量丛 $E$ 与 $F$，我们分别用 $\Gamma E$ 与 $\Gamma F$ 记 $E$ 与 $F$ 的 $C^\infty$ 截面构成的空间（确切地说，$\Gamma V$ 由满足 $\pi\circ s=\id$ 的光滑映射 $s:X\to V$ 组成，其中 $\pi:V\to X$ 是向量丛 $V$ 的投影映射）。

\begin{definition}
映射 $P:\Gamma E\to\Gamma F$ 称为\emph{线性偏微分算子}，如果它在 $X$ 的局部坐标 $x_1,\dots,x_n$ 下取形式
\[
P=\sum_{r\le m}\alpha_{i_1\cdots i_r}\,
\frac{\partial^{i_1+\cdots+i_r}}{\partial x_{i_1}\cdots\partial x_{i_r}},
\]
其中 $\alpha_{i_1\cdots i_r}(x):E_x\to F_x$ 是 $E$ 与 $F$ 的纤维之间的线性变换，光滑依赖于 $x$。
\end{definition}

换言之，$P$ 局部地是算子 $\partial/\partial x_j$ 的多项式，系数取矩阵值并光滑依赖于 $x$。形式地把未定元 $i\xi_j$ 代入 $\partial/\partial x_j$，我们得到相应的多项式 $p(x,\xi)$，其中 $\xi=(\xi_1,\dots,\xi_n)$。为说明多项式与算子之间的联系，可把算子写成 $P=p(x,D)$ 的形式，其中 $D=(-i\partial/\partial x_1,\dots,-i\partial/\partial x_n)$。\footnote{读者可能认出，映射 $\partial/\partial x_j\mapsto i\xi_j$ 就是 Fourier 变换。这一与 Fourier 分析的联系在本文后面会变得重要。} 若只保留多项式 $p(x,\xi)$ 中次数最高的那些项，我们得到齐次多项式，记作 $\sigma(x,\xi)$，称为 $P$ 的\emph{象征}。

\begin{definition}
线性偏微分算子称为\emph{椭圆的}，如果它的象征 $\sigma(x,\xi)$ 对一切 $x\in X$ 与一切非零向量 $\xi=(\xi_1,\dots,\xi_n)\neq0\in\R^n$ 都可逆。
\end{definition}

\begin{example}
设 $X=\R^n$，$E=F=\R^n\times\C$ 是平凡复线丛。此时 $\Gamma E$ 或 $\Gamma F$ 的元素就是 $\R^n$ 上的光滑复值函数。$\R^n$ 上的 Laplace 算子是二阶微分算子
\[
\nabla^2=\left(\frac{\partial}{\partial x_1}\right)^2+\cdots+\left(\frac{\partial}{\partial x_n}\right)^2 .
\]
它的象征是 $\sigma(x,\xi)=-\left((\xi_1)^2+\cdots+(\xi_n)^2\right)<0$（$\xi\neq0$），所以 Laplace 算子是椭圆的。更一般地，若取形如 $D=f\nabla^2$ 的算子，其中 $f$ 是光滑的非零复值函数，则 $D$ 也是椭圆算子。
\end{example}

\begin{example}
令 $X=\C=\R^2$，取 $z=x+iy$，其中 $x$、$y$ 是实坐标。复分析中的 Cauchy--Riemann 算子是
\[
\frac{\partial}{\partial\bar z}=\frac12\left(\frac{\partial}{\partial x}
+i\frac{\partial}{\partial y}\right).
\]
回忆 Cauchy--Riemann 微分方程由 $\partial f/\partial\bar z=0$ 给出。$\partial/\partial\bar z$ 的象征是 $\sigma(x,\xi)=\frac12 i(\xi_1+i\xi_2)$；由于 $\xi$ 是实向量，只要 $\xi\neq0$ 就有 $\sigma\neq0$，所以 $\partial/\partial\bar z$ 是椭圆算子。我们将在第 2 节末尾回到这个算子。
\end{example}

给定紧流形 $X$ 上的椭圆算子 $D:\Gamma E\to\Gamma F$，偏微分算子理论告诉我们，$D$ 的核（即 $Df=0$ 的全部解构成的空间）与 $D$ 的余核（即商空间 $\Gamma F/\operatorname{im}D$）都是有限维的。于是可以定义

% ===== 原书 p.5 =====
\begin{definition}
紧流形上椭圆算子 $D$ 的\emph{指标}由
\[
\Ind D=\dim\Ker D-\dim\Coker D
\]
给出。
\end{definition}

\begin{example}
设 $X$ 是圆周 $S^1=\R/2\pi\Z$，考虑算子 $d/d\theta$。它的象征是函数 $\sigma(x,\xi)=i\xi$，所以该算子是椭圆的。核是微分方程 $\frac{d}{d\theta}f(\theta)=0$ 的全部解构成的空间，显然就是常函数。为计算余核，注意 $S^1$ 上任何光滑函数 $f(\theta)$ 都可写成 Fourier 级数 $f(\theta)=\sum_{-\infty}^{\infty}a_ne^{in\theta}$。于是得到
\[
\frac{d}{d\theta}f(\theta)=\sum_{-\infty}^{\infty}in\,a_ne^{in\theta}.
\]
注意该级数中常数项（$n=0$）消失了，可见余核同样只是常函数。因此
\[
\Ind d/d\theta=1-1=0.
\]
\end{example}

更一般地，考虑由 $P=p(d/d\theta)$ 给出的椭圆算子，其中 $p(\xi)$ 是常系数多项式。用 Fourier 级数展开，立即可得
\[
Pf(\theta)=\sum_{-\infty}^{\infty}p(in)\,a_ne^{in\theta}.
\]
设 $n_1,\dots,n_r$ 是方程 $p(in)=0$ 的互不相同的整数解，则 $P$ 的核由一切形如 $f(\theta)=b_1e^{in_1\theta}+\cdots+b_re^{in_r\theta}$ 的函数组成，故 $\dim\Ker P=r$。又注意 $P$ 的像由一切满足 $b_{n_1}=\cdots=b_{n_r}=0$ 的函数 $f(\theta)=\sum_{-\infty}^{\infty}b_ne^{in\theta}$ 组成，于是可以把 $P$ 的余核与 $P$ 的核等同起来，故 $\dim\Coker P=r$。由此得到 $\Ind P=r-r=0$。事实上，本文后面会看到，圆周 $S^1$ 上任何线性椭圆偏微分算子的指标都是零。

按我们的定义，椭圆算子的指标是一个与某些偏微分方程的解相联系的整数。本文的目标是把指标用更容易计算的拓扑不变量表示出来，这些不变量涉及向量丛、上同调群与 Chern 类。按代数拓扑的精神，我们处理这一问题的方法可由下图很好地概括：
\[
\begin{tikzcd}[row sep=2.8em, column sep=2.8em]
K(TX) \arrow[rr, "\mathrm{t\text{-}ind}"] & & \Z \\
& \text{椭圆算子} \arrow[ul, "\sigma"'] \arrow[ur] &
\end{tikzcd}
\]

\transnote{图中三个映射分别在后面三处定义：$\sigma$ 是第 2.1 节的象征映射，$\mathrm{t\text{-}ind}$ 是第 2.2 节构造的拓扑指标，而未标出名字的对角箭头就是第 0.2 节定义的解析指标 $\Ind$（第 3.3 节）。第 2、3 节的全部内容可以概括为一句话：这个三角形交换。指标定理的证明并不是去算两边的具体数值，而是验证两条路径给出同一个整数。}


接下来三节将逐一解释图中的各个方面。我们已经看到从椭圆算子到整数的指标映射的定义。第 1 节将定义并研究函数 $K(\cdot)$ 的性质以及它与上同调的关系。在第 2 节中，我们引入 K 理论中的象征映射 $\sigma$ 与拓扑指标 $\mathrm{t\text{-}ind}$。Atiyah--Singer 指标定理的实际证明将在第 3 节给出，其内容就是证明上图交换。在这一过程中，必须把讨论扩展到椭圆伪微分算子类。

% ===== 原书 p.6 =====
\section{K 理论与上同调}

\subsection{向量丛与 $K(X)$}

设 $X$ 是紧 Hausdorff 空间，$\Vect(X)$ 是 $X$ 上复向量丛的同构类构成的空间。Whitney 和 $\oplus$ 使 $\Vect(X)$ 成为一个带单位元的阿贝尔半群。于是可以按 Grothendieck 的形式作差手续，构造相应的虚拟向量丛群 $K(X)$。

\begin{definition}
$K(X)$ 是 $\Vect(X)\times\Vect(X)$ 中 $\Delta\Vect(X)$ 的陪集构成的空间，其中 $\Delta:\Vect(X)\to\Vect(X)\times\Vect(X)$ 是对角映射。换言之，我们把 $K(X)$ 定义为向量丛的一切对 $E-F$，模去等价关系
\[
E^1-E^2=F^1-F^2
\iff
\exists\,G \text{ 使得 } E^1\oplus F^2\oplus G\cong E^2\oplus F^1\oplus G .
\]
这样得到的集合 $K(X)$ 是阿贝尔群，负元由 $-(E-F)=(F-E)$ 给出。
\end{definition}

也可以把 $K(X)$ 定义为以 $\Vect(X)$ 为生成元的自由阿贝尔群，模去由一切形如 $E+F-(E\oplus F)$ 的元素生成的子群。两种定义殊途同归。向量丛的张量积 $\otimes$ 可以延拓为 $K(X)$ 上的乘法，从而赋予 $K(X)$ 带单位元的交换环结构。

给定 $X$ 上的向量丛 $E$，用 $[E]$ 记它在 $K(X)$ 中的像。注意映射 $\Vect(X)\to K(X)$ 不必是单射。例如考虑球面 $S^2\subset\R^3$ 上的实向量丛：法丛 $NS^2$ 同构于平凡线丛 $\mathbf{1}$，于是在 $K_\R(S^2)$ 中得到
\[
[TS^2]=[3]-[NS^2]=[3]-[1]=[2],
\]
而切丛 $TS^2$ 与平凡平面丛 $[2]$ 并不同构。这正是必须在上述等价关系中引入 $G$ 项的原因。因此，若 $[E]=[F]$，则存在向量丛 $G$ 使 $E\oplus G\cong F\oplus G$；再取向量丛 $G^\perp$ 使 $G\oplus G^\perp\cong\mathbf{n}$ 对某个 $\mathbf{n}$ 成立。由此可见，$[E]=[F]$ 当且仅当 $E$ 与 $F$ 稳定等价。

\transnote{等价关系中的 $G$ 依赖于被比较的两对丛，而不是事先固定的。若把 $G$ 固定下来，「存在 $G$ 使 $E\oplus G\cong F\oplus G$」这种关系就没有传递性，无法定义陪集。取过 $G$ 后再补上正交补，一切就归约为稳定等价：$K(X)$ 的元素不必是某个具体丛的类，而是「虚拟丛」。}


给定光滑映射 $f:X\to Y$ 与 $Y$ 上的向量丛 $F$，可以构造 $X$ 上的拉回丛 $f^*F$，它对一切 $x\in X$ 满足 $(f^*F)_x=F_{f(x)}$。这一映射进而延拓为环同态 $f^*:K(Y)\to K(X)$，称为 $f$ 的诱导映射。由此可见，$K(\cdot)$ 是从紧空间范畴到带单位元交换环范畴的反变函子。特别地，若 $i:\pt\to X$ 是点到 $X$ 的包含映射，则得到映射 $i^*:K(X)\to K(\pt)\cong\Z$。\footnote{点上的向量丛就是向量空间，它由维数决定，在同构意义下唯一。} 定义\emph{约化 K 群}为 $\tilde K(X)=\Ker i^*$，就得到分裂
\[
K(X)\cong\tilde K(X)\oplus K(\pt).
\]
把 $\tilde K(X)$ 看作 $K(X)$ 模去平凡丛所得的商，则 $\tilde K(X)$ 由稳定向量丛的类组成。

\transnote{$i^*$ 就是把向量丛映为它在一点处的纤维维数，即取秩。因此 $\Ker i^*$ 由正负两部分秩相等的虚拟丛组成；$\tilde K(\pt)=0$，而 $K(\pt)=\Z$。实 K 理论中 $K_\R(\pt)=\Z$，复 K 理论中同样是 $\Z$，故上面这个分裂把「秩」这一平凡信息干净地分离出去了。}


\subsection{复形与紧支撑}

现在考虑 $X$ 不紧、而只是局部紧（且 Hausdorff）的情形。设 $X^+$ 是 $X$ 的单点紧化（若 $X$ 本身已紧，就定义 $X^+=X\cup\pt$），则可以讨论紧支撑的 K 理论。

\begin{definition}
若 $X$ 局部紧，则定义
\[
K(X):=\tilde K(X^+)=K(X^+)/K(\pt).
\]
\end{definition}

注意，若 $X$ 本身已紧，这一定义与前面的定义一致。把讨论限制在逆紧映射上，\footnote{\emph{逆紧映射}是指紧集的原像仍为紧集的映射。这一性质使逆紧映射可以延拓为两个空间的单点紧化之间的连续映射。} 这一定义就把 $K(\cdot)$ 延拓为从局部紧 Hausdorff 空间范畴到交换环范畴的反变函数（注意：$K(\cdot)$ 是带单位元的环，当且仅当 $X$ 紧）。\footnote{用紧支撑的观点，$K(\cdot)$ 也可以看作关于开集包含的共变函数。设 $i:U\to X$ 是开子集 $U$ 到 $X$ 的包含映射，定义\emph{自然延拓同态} $i_*:K(U)\to K(X)$ 为由 $X^+\to X^+/(X^+-U^+)\cong U^+$ 诱导的映射。}

% ===== 原书 p.7 =====
换一个角度，也可以考虑 $X$ 上向量丛的复形 $E^\bullet$，其形式为
\[
0\to E^0\xrightarrow{\alpha_0}E^1\xrightarrow{\alpha_1}\cdots
\xrightarrow{\alpha_{n-1}}E^n\to0,
\]
其中 $\alpha_k\alpha_{k-1}=0$。$X$ 上的两个复形 $E^\bullet$ 与 $F^\bullet$ 称为\emph{同伦的}，若存在 $X\times[0,1]$ 上的复形 $G^\bullet$，使 $E^\bullet$ 与 $F^\bullet$ 分别同构于限制 $G^\bullet|_{X\times0}$ 与 $G^\bullet|_{X\times1}$。该复形的\emph{支撑}是 $X$ 中使复形限制到 $x$ 处纤维上不成为正合列的一切点 $x$ 构成的集合。我们关心的是具有紧支撑的复形。

\begin{claim}
设 $C(X)$ 表示 $X$ 上具有紧支撑的复形按紧支撑同伦构成的同伦类集合，$C_\emptyset(X)$ 是 $C(X)$ 中由空支撑复形构成的子集，则
\[
K(X)\cong C(X)/C_\emptyset(X),
\]
其中 $K(X)$ 是上面定义的紧支撑 K 群。
\end{claim}

证明见 [6, 第 2.6 节] 或 [17, 附录]。若 $X$ 实际上是紧的，则复形 $E^\bullet$ 对应于元素
\[
\chi(E^\bullet)=\sum_k(-1)^k[E^k]\in K(X).
\]
这个交错和映射 $\chi:C(X)\to K(X)$ 就是通常的\emph{Euler 特征}。

除了任意长度的复形，也可以只用固定长度 $n>0$ 的复形，仍然得到 $K(X)\cong C^n(X)/C^n_\emptyset(X)$。特别地，把复形限制为长度 $1$ 往往很方便，即「在紧集外同构的向量丛对 $(E,F)$」。设 $\varphi:C^1(X)\to C(X)$ 是显然的包含映射，我们如下构造它的左逆 $\psi:C(X)\to C^1(X)$。给定任意长度的复形
\[
0\to E^0\xrightarrow{\alpha_0}E^1\xrightarrow{\alpha_1}\cdots
\xrightarrow{\alpha_{n-1}}E^n\to0,
\]
在每个向量丛 $E^k$ 上取 Hermite 内积 $\langle\cdot,\cdot\rangle_k$。于是得到伴随映射 $\alpha_k^*:E^{k+1}\to E^k$，由性质
\[
\langle\alpha_k(v),w\rangle_{k+1}=\langle v,\alpha_k^*(w)\rangle_k
\]
刻画。现在令 $\psi(E^\bullet)$ 为相应的长度 $1$ 复形
\[
0\to\bigoplus_i E^{2i}\xrightarrow{\alpha}\bigoplus_i E^{2i+1}\to0,
\]
其中
\[
\alpha=\bigoplus_i\alpha_{2i}+\bigoplus_i\alpha_{2i+1}^*.
\]
容易验证 $E^\bullet$ 与 $\psi(E^\bullet)$ 有相同的支撑。另外注意，复形 $\psi(E^\bullet)$ 与内积的选取无关，因为所有内积彼此同伦。

\transnote{把偶数项直接并置、奇数项先取伴随再并置，是 Hodge 理论中「$d+d^*$」这一手法的离散版本：伴随项负责把上同调「压回」核与余核。正是这个符号安排使 $\psi$ 成为左逆，从而任意长度的复形与长度 $1$ 的复形给出同一个 $K$ 群；若漏掉伴随，$\alpha$ 一般不再是复形的微分。}


% ===== 原书 p.8 =====
采用任意长度复形的好处是，乘法的定义变得方便。给定 $X$ 上的向量丛 $E$ 与 $Y$ 上的向量丛 $F$，定义它们在 $X\times Y$ 上的\emph{外张量积}为
\[
E\BT F:=\pi_X^*E\otimes\pi_Y^*F,
\]
其中 $\pi_X$ 与 $\pi_Y$ 分别是 $X\times Y$ 到 $X$ 与 $Y$ 的坐标投影。取复形的外张量积，即可把它延拓为乘积映射
\[
\BT:K(X)\otimes K(Y)\to K(X\times Y),
\]
其中 $(E^\bullet\BT F^\bullet)^i=\bigoplus_jE^j\BT F^{i-j}$，且
\[
\alpha_i^{E^\bullet\BT F^\bullet}
=\bigoplus_j\Big(\big(\alpha_j^{E^\bullet}\BT1_{F^{i-j}}\big)
+(-1)^j\big(1_{E^j}\BT\alpha_{i-j}^{F^\bullet}\big)\Big).
\]

例如，给定长度 $1$ 的复形
\[
0\to E^0\xrightarrow{\alpha}E^1\to0
\qquad\text{与}\qquad
0\to F^0\xrightarrow{\beta}F^1\to0,
\]
它们分别定义在 $X$ 与 $Y$ 上，其外张量积是 $X\times Y$ 上的复形
\[
0\to E^0\BT F^0
\xrightarrow{\alpha\BT1+1\BT\beta}
E^1\BT F^0\oplus E^0\BT F^1
\xrightarrow{-1\BT\beta+\alpha\BT1}
E^1\BT F^1\to0 .
\]
注意必须引入因子 $-1$，才能使 $\alpha_1\alpha_0=0$。按上述构造，相应的长度 $1$ 复形为
\[
0\to E^0\BT F^0\oplus E^1\BT F^1
\xrightarrow{\theta}
E^0\BT F^1\oplus E^1\BT F^0\to0,
\]
其中
\[
\theta=
\begin{pmatrix}
\alpha\BT1 & -1\BT\beta^*\\
1\BT\beta & \alpha^*\BT1
\end{pmatrix}.
\]

\subsection{齐次复形}

我们引入紧支撑 K 理论的主要动机，是考虑 $V$ 上的复形，其中 $V$ 是 $X$ 上的实向量丛，投影映射为 $\pi:V\to X$（后文将取 $V$ 为 $TX$，即 $X$ 的余切丛）。此时可以在复形上附加齐次结构。给定 $X$ 上的复向量丛 $E$ 与 $F$，可把它们提升为 $V$ 上的向量丛 $\pi^*E$ 与 $\pi^*F$，并且有 $(\pi^*E)_v=E_{\pi(v)}$ 与 $(\pi^*F)_v=F_{\pi(v)}$。若对一切 $v\in V$ 与实数 $\lambda>0$ 都有
\[
\alpha_{\lambda v}=\lambda^m\alpha_v:E_{\pi(v)}\to F_{\pi(v)},
\]
则称同态 $\alpha:\pi^*E\to\pi^*F$ 是（正）$m$ 次\emph{齐次的}。

给定 $V$ 上的度量，可见 $\alpha$ 完全由它在球丛 $S(V)$ 上的限制决定。设 $X$ 紧，$E^\bullet$ 是 $V$ 上形如
\[
0\to\pi^*E^0\xrightarrow{\alpha}\pi^*E^1\xrightarrow{\alpha}\cdots
\xrightarrow{\alpha}\pi^*E^n\to0
\]
的复形，其中 $\alpha^2=0$，且 $\alpha$ 是 $m$ 次齐次的。若 $E^\bullet$ 在 $S(V)$ 上正合，则 $E^\bullet$ 的支撑就是零截面，即 $X$ 在 $V$ 中的像。于是 $E^\bullet$ 具有紧支撑，从而代表 $K(V)$ 的一个元素。\footnote{注意，若 $m<0$，则 $\alpha$ 在零截面上不连续。所幸我们不需要 $m<0$ 的情形。不过，若 $X$ 只是局部紧的，就必须取 $m=0$，此时 $\alpha$ 在零截面上可能不连续。见 [10, 第 492--493 页]。} 我们断言，$K(V)$ 可以用这类齐次复形来定义。

\transnote{原文用 $TX$ 记 $X$ 的\emph{余切}丛，而通行记号 $TX$ 表示切丛，余切丛写作 $T^*X$。本译本保留原文字面，因此下文凡出现 $K(TX)$，都应理解为余切丛的 K 群。这样记也有好处：象征 $\sigma(x,\xi)$ 是余向量 $\xi$ 的齐次多项式，其取值自然落在余切丛的 K 群中。}

% ===== 原书 p.9 =====
\begin{claim}
记 ${}^{m}C(V)$ 为 $V$ 上 $m$ 次齐次复形的同伦类集合，${}^{m}C_\emptyset(V)$ 为其中使同态 $\alpha$ 沿球丛 $S(V)$ 的每个纤维为常值的那部分正合复形。则
\[
{}^{m}C(V)/{}^{m}C_\emptyset(V)\cong C(V)/C_\emptyset(V)\cong K(V).
\]
\footnote{与前述相同，可以把讨论限制在固定长度的齐次复形上。}
\end{claim}

\begin{proof}
考虑包含 ${}^{m}C(V)\subset C(V)$，我们想构造一个逆映射 $C(V)\to{}^{m}C(V)$。设 $E^\bullet\in C(V)$ 是支撑为紧集 $L$ 的复形
\[
0\longrightarrow E^0\xrightarrow{\ \alpha_0\ }E^1\xrightarrow{\ \alpha_1\ }\cdots\xrightarrow{\ \alpha_{n-1}\ }E^n\longrightarrow 0,
\]
在 $V$ 上取一个度量，并令 $D(V)$ 是包含 $L$ 的圆盘丛。把 $E^i$ 限制到零截面，得 $F^i=E^i|_X$，于是在 $D(V)$ 上有 $E^i\cong\pi^*F^i$。定义 $\beta_i:\pi^*F^i\to\pi^*F^{i+1}$ 为在相应的球丛 $S(V)$ 上与 $\alpha_i$ 一致的那个 $m$ 次齐次映射，这样就得到 $V$ 上的一个齐次复形 $F^\bullet$：
\[
0\longrightarrow\pi^*F^0\xrightarrow{\ \beta_0\ }\pi^*F^1\xrightarrow{\ \beta_1\ }\cdots\xrightarrow{\ \beta_{n-1}\ }\pi^*F^n\longrightarrow 0.
\]
容易看出 $E^\bullet$ 与 $F^\bullet$ 同伦，因此它们对应 $K(V)$ 中的同一个元素；于是映射 $C(V)\to{}^{m}C(V)$ 是同构。其次，若 $E^\bullet\in C_\emptyset(V)$ 的支撑为空，则 $E^\bullet$ 同伦于一个使 $\alpha_i$ 沿 $V$ 的每个纤维为常值的复形，由此可见 $F^\bullet\in{}^{m}C_\emptyset(V)$。反之，若 $F^\bullet\in{}^{m}C_\emptyset(V)$ 使 $\beta_i$ 沿 $S(V)$ 的每个纤维为常值，则取同伦 $\alpha_i(v,t)=\|v\|^{tm}\beta_i(v/\|v\|)$，可见 $F^\bullet$ 同伦于某个正合复形 $E^\bullet\in C_\emptyset(V)$。于是 ${}^{m}C(V)/{}^{m}C_\emptyset(V)\cong C(V)/C_\emptyset(V)$。
\end{proof}

\subsection{Thom 同构}
% ===== 原书 p.9 =====
设$V$是紧空间$X$上的复向量丛。外代数（exterior algebra）$\Lambda^*(V)$\footnote{第$n$个外幂$\Lambda^n(V)$是$n$重张量积$V^{\otimes n}$的斜对称化。}于是给出$V$上一个次数为 1 的、由向量丛组成的齐次复形$\Lambda(V)$，称为外复形（exterior complex）：
\[
0\longrightarrow\pi^*\Lambda^0(V)\xrightarrow{\ \alpha\ }\pi^*\Lambda^1(V)\xrightarrow{\ \alpha\ }\cdots\xrightarrow{\ \alpha\ }\pi^*\Lambda^n(V)\longrightarrow 0,
\]
其中$\pi:V\to X$，且对一切$v\in V$与$w\in\pi^*\Lambda^k(V)_v=\Lambda^k(V)_{\pi(v)}$，有$\alpha:(v,w)\mapsto(v,v\wedge w)$。由于$\alpha^2=0$，且该复形在零截面之外正合，可见它定义一个元素$\lambda_V\in K(V)=\tilde K(V^+)$；这里$V^+=X^V=D(V)/S(V)=P(V\oplus\mathbf{1})/P(V)$是$V$的 Thom 空间\footnote{这里把 Thom 空间定义为全空间的单点紧化。若$X$只是局部紧，则须改为逐个纤维取单点紧化。}。下面陈述$K$理论的基本定理。

\begin{theorem}[Thom 同构定理]
若$V$是紧空间$X$上的复向量丛，且把$K(V)$通过提升$\pi^*:K(X)\to K(V)$看作$K(X)$-模，则由乘以$\lambda_V$给出的 Thom 映射$\varphi:K(X)\to K(V)$是同构。
\end{theorem}

该定理的证明在附录中给出梗概。若$X$只是局部紧，则复形$\Lambda(V)$没有紧支撑。然而，若$E^\bullet$是$V$上具有紧支撑的复形，则$\Lambda(V)$与$E^\bullet$的乘积确有紧支撑，Thom 同构定理仍然成立。设$s:X\to V$是零截面，注意当$X$紧时
\[
s^*\lambda_V=s^*\circ\varphi(1)=\sum_{k=0}^n(-1)^k\Lambda^k(V).
\]
\transnote{$\lambda_V$落在$K(V)$里，而$K(V)$是$V$上紧支撑$K$理论，它靠单点紧化才等同$\tilde K(V^+)$。所以零截面上的$s^*\lambda_V$是一个$K(X)$中的类，不能把$\Lambda^k(V)$直接当作$X$上的丛来读；须先沿$s$拉回。}

% ===== 原书 p.10 =====
注意$\Lambda^*(V\oplus W)\cong\Lambda^*(V)\otimes\Lambda^*(W)$，可见$\lambda_{V\oplus W}=\lambda_V\cdot\lambda_W$，从而 Thom 同构是传递的。换言之，同构$K(X)\to K(V\oplus W)$与复合$K(X)\to K(V)\to K(V\oplus W)$相同。

\begin{example}
若$X$是一点且$V=\C^n$，则得到 Thom 同构
\[
\varphi:K(\pt)\xrightarrow{\ \cong\ }K(\C^n)\cong\tilde K(S^{2n}),
\]
从而$\tilde K(S^{2n})$同构于$\Z$，并由 Bott 类（Bott class）$\lambda_{\C^n}=\lambda_n$生成。于是$K(S^{2n})\cong\Z^2$。
\end{example}

\subsection{一个周期的上同调理论}
$K$理论最有趣也最有用的特点之一，是它具有广义上同调理论的特征性质。为看出这一点，先需要几个定义。

\begin{definition}
带基点空间$X$的约化悬挂（reduced suspension）由
\[
SX=S^1\wedge X=(S^1\times X)/(S^1\vee X)
\]
给出；$X$的$n$次迭代悬挂由
\[
S^nX=\underbrace{SS\cdots S}_{n}X=S^n\wedge X
\]
给出，其中$X\vee Y=(X\times\pt)\cup(\pt\times Y)$是两个带基点空间的单点并（one point union），$X\wedge Y=X\times Y/X\vee Y$是 smash 积。
\end{definition}

\begin{definition}
对$n\ge 0$，用
\begin{align*}
K^{-n}(X,Y)&:=\tilde K(S^n(X/Y)),\\
K^{-n}(X)&:=K^{-n}(X,\varnothing)=\tilde K(S^nX^+)
\end{align*}
定义次数平移（degree shift）\footnote{用紧支撑$K$理论，等价的定义是$K^{-n}(X):=K(\R^n\times X)$。}。
\end{definition}

\begin{corollary}[Bott 周期性]
由
\[
K^{-n}(X)\xrightarrow{\ \lambda_1\BT\cdot\ }K^{-n-2}(X)
\]
给出的映射对一切$n\ge 0$都是同构，其中$\lambda_1$是生成$\tilde K(S^2)=K^{-2}(\pt)$的 Bott 类。
\end{corollary}

\begin{proof}
取$V=\C^m\times X$为平凡丛，得到 Thom 同构
\begin{align*}
K(X)&\xrightarrow{\ \cong\ }K(\C^m\times X)=\tilde K\big((\C^m\times X)^+\big)\\
&=\tilde K(S^{2m}\wedge X^+)=K^{-2m}(X),
\end{align*}
它由
\[
x\mapsto\lambda_{X\times\C^m}\cdot\pi^*x=\pi^*_{\C^m}\lambda_m\cdot\pi^*_Xx=\lambda_m\BT x
\]
给出；这里把外张量积（external tensor product）看作映射$\BT:\tilde K(X)\otimes\tilde K(Y)\to\tilde K(X\wedge Y)$。注意$\lambda_m=(\lambda_1)^m$，所要结论立得。
\end{proof}

\transnote{此处的外张量积是带帽的$\widehat{\otimes}$，不是普通张量积。它把两个约化$K$群送进$X\wedge Y$的约化$K$群，正是悬挂同构在$K$理论里的化身；Bott 类$\lambda_m$按$\BT$相乘时才有$\lambda_m=(\lambda_1)^m$。}

% ===== 原书 p.11 =====
由 Bott 周期性，可见所有偶数维群都同构于$K^0(X)=K(X)$，所有奇数维群都同构于$K^1(X)=\tilde K(SX^+)$。于是得到$\Z_2$-分次环$K^*(X)=K^0(X)\oplus K^1(X)$，它给出一个周期的上同调理论，并伴有如下正合的「环」序列：
\[
\begin{tikzcd}[column sep=3.2em,row sep=2.4em]
K^0(X,Y)\arrow[r] & K^0(X)\arrow[r] & K^0(Y)\arrow[d]\\
K^1(Y)\arrow[u] & K^1(X)\arrow[l] & K^1(X,Y)\arrow[l]
\end{tikzcd}
\]
利用这个正合列，可以得到与整上同调类似的结果。例如，下面命题的证明留给读者（或见[6, §2.5.2]）。

\begin{proposition}
若$X$是有限 CW 复形，且胞腔只出现在偶数维，则$K^0(X)\cong\Z^r$，其中$r$是胞腔数，且$K^1(X)=0$。于是$K^*(X)\cong\Z^r$。
\end{proposition}

于是，回想复射影空间$\C\PP^n$在$0,2,\ldots,2n$每一维恰有一个胞腔，可见$K^*(\C\PP^n)\cong\Z^{n+1}$。仍得到$K(S^{2n})\cong\Z^2$。下面给出$K$理论中 K\"unneth 公式的类似物（见[6, §2.7.15]）。

\begin{theorem}[K\"unneth 定理]
设$X$、$Y$是有限 CW 复形。则有自然正合列（指标取值于$\Z_2$）
\[
0\to\sum_{i+j=k}K^i(X)\otimes K^j(Y)\to K^k(X\times Y)\to\sum_{i+j=k+1}\operatorname{Tor}\big(K^i(X),K^j(Y)\big)\to 0.
\]
\end{theorem}

\begin{example}
由于圆$S^1$上所有复向量丛都平凡，有$K(S^1)\cong\Z$。于是对一点的情形得到$K^0(\pt)\cong\Z$与$K^1(\pt)=\tilde K(S^1)=0$。另一方面，对圆$S^1$，可见$K^0(S^1)\cong\Z$与$K^1(S^1)=\tilde K(S^2)\cong\Z$。于是由 K\"unneth 定理得到$K(X\times S^1)\cong K^0(X)\oplus K^1(X)=K^*(X)$\footnote{也可取$K(X\times S^1)\cong K(X\wedge S^1)\oplus K(X)\oplus K(S^1)\oplus K(\pt)\cong\tilde K(SX^+)\oplus K(X)$。}。
\end{example}

\subsection{Chern 特征}
为讨论$K$理论与上同调的关系，需要引入示性类的概念。这里使用 \v{C}ech 上同调；处理局部紧空间时取紧支撑上同调。给定底空间$X$上的复$n$-平面丛$E$，对每个$i=0,\ldots,n$为它配以 Chern 类$c_i(E)\in H^{2i}(X;\Z)$，其中$c_0(E)=1$；并把总 Chern 类（total Chern class）定义为形式和
\[
c(E)=1+c_1(E)+\cdots+c_n(E).
\]
Chern 类的构造见[16, §14]或[13, §20]。这里只需知道 Chern 类满足下面四条重要性质\footnote{为排除所有 Chern 类都为零的退化情形，还要求$c_1(\text{Hopf 丛})$是$H^2(\C\PP^n)$的典范生成元（$n\ge 1$）。这一事实连同自然性性质与乘积公式，足以完全刻画 Chern 类。}：

\noindent\textbf{性质（自然性）}。若存在复$n$-平面丛之间的丛映射$E\to F$覆盖$f:X\to Y$，其中$E$、$F$分别是$X$、$Y$上的复$n$-平面丛，则$c(E)=f^*c(F)$。

% ===== 原书 p.12 =====
\noindent\textbf{性质（乘积公式）}。若$E$、$F$分别是同一底空间$X$上的复$m$-平面丛与复$n$-平面丛，则总 Chern 类满足
\[
c(E\oplus F)=c(E)\smile c(F),\quad\text{即}\quad c_i(E\oplus F)=\sum_{k=0}^{i}c_k(E)\smile c_{i-k}(F)\quad(i\ge 0).
\]

\noindent\textbf{性质}。若把复向量丛$E$的共轭丛（conjugate bundle）记为$\overline{E}$，于是$\overline{E}$与$E$的底实空间相同、复结构相反，则$c_k(\overline{E})=(-1)^kc_k(E)$；用总 Chern 类表示即
\[
c(\overline{E})=1-c_1(E)+c_2(E)-\cdots\pm c_n(E).
\]

\noindent\textbf{性质}。若$\mathbf{n}$是平凡复$n$-平面丛，则$c(\mathbf{n})=1$。

设$X$上的复向量丛$E$分裂为复线丛的直和$E=L_1\oplus\cdots L_n$。令$x_i=c_1(L_i)\in H^2(X;\Z)$，乘积公式给出
\[
c(E)=\prod_i c(L_i)=\prod_i(1+x_i)=1+\sum_i x_i+\sum_{i<j}x_ix_j+\cdots,
\]
可见 Chern 类$c_k(E)$是诸$x_i$的第$k$个初等对称函数（elementary symmetric function）。于是$x_i$的任一对称函数都可表为 Chern 类的多项式。虽然并非每个向量丛都这样分解\footnote{不分裂的最简单例子是$T\C\PP^2$。}，但有下面这条重要定理。

\begin{theorem}[分裂原理]
若$E$是$X$上的复向量丛，则存在空间$F(E)$与映射$\pi:F(E)\to X$，使$\pi^*:H^*(X)\to H^*(F(E))$是单射，且$\pi^*E$分裂为复线丛的直和$E=L_1\oplus\cdots\oplus L_n$。
\end{theorem}

更确切地说，取$F(E)$为$E$的旗丛（flag bundle）。证明见[6, §2.7.1]或[13, §21]。分裂原理与 Chern 类的自然性结合，说明在只涉及 Chern 类的计算中，任何向量丛都可当作线丛的直和。特别地，可以定义：

\begin{definition}
复$n$-平面丛$E$的 Chern 特征由下式给出：
\begin{align*}
\ch(E)&=\sum_i e^{x_i}=n+\sum_i x_i+\frac{1}{2!}\sum_i x_i^2+\cdots\\
&=n+c_1+\frac12(c_1^2-2c_2)+\cdots
\end{align*}
\end{definition}

\transnote{分裂原理的用法是「先拉到$F(E)$上算，再拉回」。$\pi^*$单射保证$H^*(X)$里的等式可在$H^*(F(E))$中检验，所以只用 Chern 类写出的示性类，都能在$\pi^*E$按线丛算完后唯一确定。$\ch$、$\Td$都是这样定义的。}

对线丛$L_1$、$L_2$，注意$c_1(L_1\otimes L_2)=c_1(L_1)+c_1(L_2)$，于是
\[
\ch(L_1\otimes L_2)=e^{x_1+x_2}=e^{x_1}e^{x_2}=\ch(L_1)\smile\ch(L_2).
\]
由此对任意复向量丛$E$、$F$得到
\begin{align*}
\ch(E\oplus F)&=\ch(E)+\ch(F),\\
\ch(E\otimes F)&=\ch(E)\smile\ch(F).
\end{align*}
把 Chern 特征扩充到$K(X)$，便得到自然环同态
\[
\ch:K(X)\to H^{\mathrm{even}}(X;\Q)
\]
% ===== 原书 p.13 =====
即从$K(X)$到$X$的偶数维有理上同调的自然环同态。进一步，把$\ch:\tilde K(SX)\to\tilde H^{\mathrm{even}}(SX^+;\Q)$与$\alpha^{-1}$复合，其中$\alpha:H^{\mathrm{odd}}(X;\Q)\to\tilde H^{\mathrm{even}}(SX^+;\Q)$是乘以$H^1(S^1;\Z)$的典范生成元给出的悬挂同构，就能把 Chern 特征再扩充，得到自然环同态
\[
\ch:K^1(X)\to H^{\mathrm{odd}}(X;\Q).
\]
于是可把 Chern 特征看作函子$K^*(\cdot)\to H^*(\cdot;\Q)$的自然变换；若用有理$K$理论消去挠，就得到

\begin{theorem}
若$X$是有限 CW 复形，则 Chern 特征
\[
\ch:K^*(X)\otimes\Q\to H^*(X;\Q)
\]
是有理$K$理论与有理上同调之间的自然同构。
\end{theorem}

谱序列证明见[8, §2]。偶维球面的特例在[15, §3]中结合 Bott 周期性定理的证明讨论。

\subsection{上同调的 Thom 同构}
其次，设$E$是$X$上的定向实$n$-平面丛。已知存在唯一生成元$u\in H^n(E,E_0;\Z)=H^n_c(E;\Z)$，其中$H^*_c(\cdot)$是紧支撑上同调，$E_0$是去掉$E$的零截面所得的删去空间（deleted space），使得$u$在每个纤维上的限制诱导预先给定的定向。这个生成元$u$称为$E$的上同调 Thom 类（cohomology Thom class）。再设$i:X\to E$是零截面，用$u$在$X$上的限制定义$E$的 Euler 类：
\[
e(X)=i^*u\in H^n(X;\Z).
\]
若$E$是$X$上的复$n$-平面丛，则$E$的复结构在其底实$2n$-平面丛$E_\R$上诱导一个定向。于是最高 Chern 类由$c_n(E)=e(E_\R)$定义。下面陈述上同调 Thom 同构定理。

\begin{theorem}[Thom 同构]
由$\psi:x\mapsto\pi^*x\smile u$给出的 Thom 同态
\[
\psi:H^k(X)\to H^{n+k}_c(E)
\]
是同构，其中$\pi:E\to X$是投影\footnote{与$K$理论不同，上同调的 Thom 同构定理在$E$为实丛时也适用。}。
\end{theorem}

\transnote{原文把$E$的 Euler 类记作$e(X)=i^*u$，而不是$e(E)$。原因是零截面把$X$嵌入$E$，拉回后类落在$H^n(X;\Z)$上；到第 2.3 节写$e(X)=e(TX)$时沿用同一写法。读时不要误以为这里算的是底流形的 Euler 类。}

证明见[16, §10]。用 Chern 特征可以比较$K$理论与上同调的 Thom 同构。设$E$是紧空间$X$上的复$n$-平面丛，得到下图：
\[
\begin{tikzcd}[column sep=4.2em,row sep=2.6em]
K(X)\arrow[r,"\varphi"]\arrow[d,"\ch"] & K(E)\arrow[d,"\ch"]\\
H^*(X;\Q)\arrow[r,"\psi"] & H^*_c(E;\Q)
\end{tikzcd}
\]
其中竖直映射是相应的 Chern 特征，水平映射是 Thom 同构$\varphi:x\mapsto\pi^*x\cdot\lambda_E$与$\psi:x\mapsto\pi^*x\smile u$。由于该图不交换，需要引入修正因子。确切地说，令
\[
\psi\big(\mu(E)\smile\ch(x)\big)=\ch\big(\varphi(x)\big),
\]
% ===== 原书 p.14 =====
其中
\[
\mu(E)=\psi^{-1}\circ\ch\circ\varphi(1)=\psi^{-1}\ch(\lambda_E)\in H^*(X;\Q).
\]
下面把这个上同调类$\mu(E)$明确算出。

设$i:X\to E$是零截面，回想$i^*\lambda_E=\sum_i(-1)^i\Lambda^i(E)$，并注意对一切上同调类$x\in H^*(X;\Q)$有$i^*\psi(x)=x\smile e(E_\R)$。特别地，取$x=\mu(E)$，由 Chern 类的自然性得到
\begin{align*}
\mu(E)\smile e(E_\R)&=i^*\psi\big(\mu(E)\big)=i^*\ch(\lambda_E)\\
&=\ch(i^*\lambda_E)=\ch\Big(\sum_{i=0}^n(-1)^i\Lambda^i(E)\Big).
\end{align*}
于是援引分裂原理，可设$E$分解为直和$E=L_1\oplus\cdots\oplus L_n$。注意$\Lambda^i(E)\cong\bigoplus_{j_1<\cdots<j_i}L_{j_1}\otimes\cdots\otimes L_{j_i}$，得到
\begin{align*}
\mu(E)\smile e(E_\R)&=\ch\Big(\sum_{i=0}^n(-1)^i\Lambda^i(E)\Big)\\
&=\ch\Big(\prod_{i=1}^n(1-L_i)\Big)=\prod_{i=1}^n(1-e^{x_i}),
\end{align*}
其中$x_i=c_1(L_i)\in H^2(X;\Z)$，$i=1,\ldots,n$。由$e(E_\R)=c_n(E)=\prod_i x_i$得出
\[
\mu(E)=\prod_{i=1}^n\frac{1-e^{x_i}}{x_i}.
\]
为进一步化简，再引入一个示性类。

\begin{definition}
复向量丛$E$的 Todd 类由下式给出：
\[
\Td(E)=\prod_{i=1}^n\frac{x_i}{1-e^{-x_i}}=1+\frac12c_1+\frac{1}{12}\big(c_2+c_1^2\big)+\cdots
\]
\end{definition}

注意由定义，Todd 类对任意两个向量丛$E$、$F$满足$\Td(E\oplus F)=\Td(E)\smile\Td(F)$，且对平凡复$n$-平面丛$\mathbf{n}$有$\Td(\mathbf{n})=1$。于是可见，用 Todd 类表示，修正因子为
\[
\mu(E)=(-1)^n\Td(\overline{E})^{-1};
\]
这里用到 Todd 类可逆，因其常数项为 1。

\transnote{定义式的分母是$1-e^{-x_i}$，不是$1-e^{x_i}$；换成后者会得到$\Td(E)^{-1}$那一族。修正因子里$\overline{E}$的共轭与$(-1)^n$缺一不可：$\Td(\overline{E})^{-1}=\prod_i(e^{x_i}-1)/x_i$，乘上$(-1)^n$才等于$\prod_i(1-e^{x_i})/x_i$。}

\section{指标定理的陈述}
% ===== 原书 p.14 =====
\subsection{象征映射与椭圆复形}
设紧空间$X$上有椭圆算子$D:\Gamma E\to\Gamma F$。于是对每个$x\in X$与$\xi\in TX_x$，$D$的象征$\sigma(x,\xi)$是同态$E_x\to F_x$，它随$x$、$\xi$光滑变化。可见这一族同态给出一个齐次复形
\[
0\longrightarrow\pi^*E\xrightarrow{\ \alpha\ }\pi^*F\longrightarrow 0,
\]
其底空间为$TX$。\footnote{本节与下一节用$TX$表示$X$的余切丛（cotangent bundle）。注意若在$X$上取定度量，则可把$X$的切丛与余切丛等同。}由于$D$椭圆，已知$\xi\neq 0$时$\sigma(x,\xi)$可逆，故该复形在零截面之外正合。可见这个复形定义$K(TX)$中的一个元素，称为$D$的象征，记作$\sigma(D)$。

% ===== 原书 p.15 =====
这一构造几乎不加改动即可推广到微分算子复形。给定紧空间$X$上的一列向量丛$E^0,\ldots,E^n$与偏微分算子$d_i:\Gamma E^i\to\Gamma E^{i+1}$，称复形
\[
0\longrightarrow\Gamma E^0\xrightarrow{\ d_0\ }\Gamma E^1\xrightarrow{\ d_1\ }\cdots\xrightarrow{\ d_{n-1}\ }\Gamma E^n\longrightarrow 0
\]
为椭圆复形，若$d_{i+1}d_i=0$，且相应的象征复形
\[
0\longrightarrow\pi^*E^0\xrightarrow{\ \sigma_0\ }\pi^*E^1\xrightarrow{\ \sigma_1\ }\cdots\xrightarrow{\ \sigma_{n-1}\ }\pi^*E^n\longrightarrow 0
\]
在$TX$上于零截面之外正合。与前面一样，这个复形给出一个元素$\sigma(D^\bullet)\in K(TX)$，称为椭圆复形$D^\bullet$的象征。对椭圆复形，取代指标的是欧拉示性数（Euler characteristic），它由
\[
\chi(D^\bullet)=\sum_{i=0}^n(-1)^i\dim H^i(D^\bullet)
\]
定义，其中$H^i(D^\bullet)=\Ker d_i/\operatorname{Im} d_{i-1}$是复形的上同调。注意当复形$D^\bullet$只是单个算子时，欧拉示性数化为前面关于指标的定义。进一步，若按§1构造相应的长度为 1 的复形，便得到椭圆算子
\[
D:\Gamma\Big(\bigoplus_iE^{2i}\Big)\to\Gamma\Big(\bigoplus_iE^{2i+1}\Big),
\]
它由$D=\bigoplus_i d_{2i}+\bigoplus_i d^*_{2i+1}$给出，其中$d^*$表示关于某个度量的伴随。于是注意$\sigma(D)=\sigma(D^\bullet)\in K(TX)$且$\Ind D=\chi(D^\bullet)$，问题就化归为单个椭圆算子的情形。

\begin{example}
考虑（复化后的）de Rham 复形
\[
0\longrightarrow\Gamma\Lambda^0(TX\otimes\C)\xrightarrow{\ d\ }\Gamma\Lambda^1(TX\otimes\C)\xrightarrow{\ d\ }\cdots\xrightarrow{\ d\ }\Gamma\Lambda^n(TX\otimes\C)\longrightarrow 0,
\]
其中$\Gamma\Lambda^i(TX)$是$X$上次数为$i$的微分形式空间，$d$是外微分（exterior derivative）。此时复形的上同调就是 de Rham 上同调，$H^i(D^\bullet)=H^i_{\mathrm{dR}}(X)\otimes\C$，于是 de Rham 复形的欧拉示性数等于流形$X$的欧拉示性数：
\[
\chi(D^\bullet)=\sum_i\dim H^i_{\mathrm{dR}}(X)=\sum_i(-1)^i\dim H^i(X;\R)=\chi(X).
\]
相应的元素$\sigma(D^\bullet)\in K(TX)$称为 de Rham 象征，记作$\rho_X$。
\end{example}

\transnote{原文第一项写作$\sum_i\dim H^i_{\mathrm{dR}}(X)$，缺少$(-1)^i$。按欧拉示性数的定义应为$\sum_i(-1)^i\dim H^i_{\mathrm{dR}}(X)$；此处照译保留，读时须补符号。}

\subsection{拓扑指标的构造}
由于象征给出从$X$上椭圆算子到$K(TX)$的映射，我们希望把它与适当的映射$K(TX)\to\Z$复合。最先想到的一个，是点包含$i:\pt\to TX$诱导的映射。然而，回想定义$K(TX)=\tilde K(TX^+)$，会失望地看到映射$i^*:K(TX)\to K(\pt)\cong\Z$恒为零。事实上，注意$i$延拓为映射$i:\pt\to TX^+$，而约化$K$群$\tilde K(TX^+)$按定义是$\Ker i^*$，上述结论立即得到。所以必须到更深处寻找所需的映射。

% ===== 原书 p.16 =====
与其考虑$X$中基点的包含，不如整体地思考。回想若$X$是紧流形，则可把$X$嵌入适当大的$\R^n$\footnote{这可用涉及单位分解的简单构造说明。}。给定这样的嵌入$X\subset\R^n$，可把它延拓为切丛的嵌入$TX\subset T\R^n$。由 Thom 同构定理$K(T\R^n)\cong\Z$，故只要构造出映射$K(TX)\to K(T\R^n)$就满意了。

更一般地，若$i:X\to Y$是紧子流形$X$到$Y$的包含，可把它延拓为切丛的包含$i:TX\to TY$。随后如下构造映射
\[
i_!:K(TX)\to K(TY)
\]
（注意该映射关于包含是函子性的）。取定$Y$上的度量，令$N$是$X$在$Y$中的管状邻域，它与$X$在$Y$中的法丛微分同胚。延拓到切丛，可见$TN$也是$TX$在$TY$中的管状邻域，与$TX$在$TY$中的法丛微分同胚。管状邻域定理的证明见[16, §11]。由于$TN$是$TX$上的实向量丛，我们希望给它一个复结构，以便应用 Thom 同构。

考虑$X$上任一向量丛$E$，回想$E$局部同胚于$U\times\R^m$，其中$U\subset X$同胚于$\R^n$的一个开区域。于是$E$在$(x,\xi)\in E$处的切空间为$TE_{(x,\xi)}=TU_x\times T\R^m_\xi\cong TX_x\oplus E_x$，可见$E$上的切丛$TE$有分解
\[
TE=\pi^*TX\oplus\pi^*E,
\]
其中$\pi:E\to X$是投影。特别地，取$E=TX$与$E=TY|_X$，得到
\begin{align*}
T(TX)&\cong\pi^*TX\oplus\pi^*TX,\\
T(TY)|_{TX}&\cong\pi^*(TY|_X)\oplus\pi^*(TY|_X),
\end{align*}
其中$T(TX)$与$T(TY)|_{TX}$都是$TX$上的丛，$\pi:TX\to X$是投影。又因$N$与$TN$分别是$X$在$Y$中、$TX$在$TY$中的法丛，有$N\oplus TX\cong TY|_X$与$TN\oplus T(TX)\cong T(TY)|_{TX}$。于是$TN$分解为
\[
TN\cong\pi^*N\oplus\pi^*N.
\]
这样就能给$TN$赋予复向量丛的结构：
\[
TN\cong\pi^*N\oplus i\pi^*N\cong(\pi^*N)\otimes_\R\C\cong\pi^*(N\otimes_\R\C).
\]
于是把 Thom 同构$\varphi:K(TX)\to K(TN)$与自然的延拓同态$h:K(TN)\to K(TY)$复合，便得到所需映射
\[
i_!:K(TX)\xrightarrow{\ \varphi\ }K(TN)\xrightarrow{\ h\ }K(TY).
\]
注意该映射不依赖邻域$N$的选取，且是函子性的（即$(j\circ i)_!=j_!\circ i_!$），因为 Thom 同构传递。还有
\[
i^*\circ i_!:x\mapsto i^*\lambda_{TN}\cdot x=\Big(\sum_{i=0}^n(-1)^i\Lambda^i(N\otimes_\R\C)\Big)\cdot x.
\]

% ===== 原书 p.17 =====
\begin{definition}
若$X$是紧流形，取充分大的$m$与嵌入$i:X\to\R^m$，并令$j:P\to\R^m$是原点的包含。注意$j_!:K(TP)\to K(T\R^m)$就是 Thom 同构$\varphi:K(\pt)\to K(\C^m)$，故可把拓扑指标定义为复合
\[
\tind:K(TX)\xrightarrow{\ i_!\ }K(T\R^m)\xrightarrow{\ j_!^{-1}\ }K(TP)\cong\Z.
\]
\end{definition}

设已有两个嵌入$i:X\to\R^m$与$i':X\to\R^n$。考虑对角嵌入$i\oplus i':X\to\R^{m+n}$，可见$i\oplus i'$同痕（isotopic）于嵌入$i\oplus 0$与$0\oplus i'$。令$\lambda_n\in K(\C^n)$表示$\C^n\cong T\R^n$的 Thom 类，注意$(i\oplus 0)_!(x)=i_!(x)\cdot\lambda_n$与$(0\oplus i')_!(x)=\lambda_m\cdot i'_!(x)$。于是由 Thom 同构的传递性，$i\oplus 0$与$i$确定同一个 t-ind，$0\oplus i'$与$i'$亦然。可见拓扑指标不依赖嵌入的选取。

现已构造出所需映射$\tind:K(TX)\to\Z$，可把它与象征映射$\sigma:(\text{椭圆算子})\to K(TX)$复合。终于可以陈述本文的主要结果。

\begin{theorem}[Atiyah–Singer 指标定理]
若$D$是紧流形上的椭圆算子，则
\[
\Ind D=\tind\sigma(D).
\]
\end{theorem}

\transnote{$i^*:K(TX)\to K(\pt)$恒为零，是因为$\tilde K(TX^+)$就定义为$\Ker i^*$。所以指标不能靠「在一点处取纤维」读出；第 2.2 节绕开的正路是嵌入$\R^m$、把法丛复化，再用 Thom 同构造出$\tind$。}

\subsection{指标定理的上同调形式}
在进入 Atiyah–Singer 指标定理的$K$理论证明之前，先短暂离题，用示性类与上同调重新表述该定理。1963 年 Atiyah 与 Singer 最初证明时用的正是这种形式，工具来自配边理论。$K$理论表述或许能给出更优雅、更一般的证明，而上同调形式便于直接计算例子。为用上同调写出拓扑指标，设$X$是嵌入$\R^k$的紧$n$维流形，考虑如下非交换图：
\[
\begin{tikzcd}[column sep=3.4em,row sep=2.8em]
K(TX)\arrow[r,"\varphi"]\arrow[d] & K(TN)\arrow[r,"h"]\arrow[d] & K(T\R^k)\arrow[d] & K(TP)\cong\Z\arrow[l,"\varphi'"']\arrow[d]\\
H^*_c(TX)\arrow[r,"\psi"] & H^*_c(TN)\arrow[r,"k"] & H^*_c(T\R^k) & H^*_c(TP)\cong\Q\arrow[l,"\psi'"']
\end{tikzcd}
\]
其中$\varphi$、$\psi$、$\varphi'$、$\psi'$是各种 Thom 同构，$h$与$k$是延拓同态，竖直映射是 Chern 特征$\ch:K(\cdot)\to H^{\mathrm{even}}_c(\cdot;\Q)$。此图虽然不交换，但回想§1：对$X$上的复向量丛$E$，$K$理论与上同调的 Thom 同构由下式相联系：
\[
\psi^{-1}\circ\ch\circ\varphi:x\mapsto(-1)^n\Td(\overline{E})^{-1}\smile\ch(x).
\]

此处因$\overline{T\R^k}$显然是点$TP$上的平凡丛，有$\Td(\overline{T\R^k})=1$。又因$TN\cong\pi^*N\otimes_\R\C$，可见$TN\cong\overline{TN}$。注意$T(TX)\oplus TN=T(T\R^k)|_{TX}$是$TX$上的平凡丛，且$T(TX)\cong\pi^*TX\oplus\pi^*TX\cong\pi^*TX\otimes_\R\C$，得到
\[
\Td(\overline{TN})^{-1}=\Td(TN)^{-1}=\Td(T(TX))=\pi^*\Td(TX\otimes_\R\C)=\pi^*\IndexClass(X),
\]
% ===== 原书 p.18 =====
其中最后一个等式给出流形$X$的指标类$\IndexClass(X)$的定义。于是类$\Td(\overline{TN})^{-1}=\pi^*\IndexClass(X)$与嵌入无关。因此
\begin{align*}
\psi'^{-1}\circ\ch\circ\varphi'&:x\mapsto(-1)^k\Td(\overline{T\R^k})^{-1}\smile\ch(x)=(-1)^k\ch(x),\\
\psi^{-1}\circ\ch\circ\varphi&:x\mapsto(-1)^{n-k}\Td(\overline{TN})^{-1}\smile\ch(x)=(-1)^{n-k}\IndexClass(X)\smile\ch(x).
\end{align*}
回想$\tind=\varphi'^{-1}\circ h\circ\varphi$，经一番扩展的图追踪得到
\[
\ch\circ h\circ\varphi(x)=\ch\circ\varphi'(\tind x)=(-1)^k\psi'(\tind x)=(-1)^k(\tind x)\psi'(1),
\]
其中$\psi'(1)$是$H^*_c(T\R^k)$的典范生成元。再令$[T\R^k]$为$T\R^k$的基本同调类（fundamental homology class），有$\psi'(1)[T\R^n]=1$，于是算出
\begin{align*}
\tind x&=(\tind x\smile\psi'(1))[T\R^n]=(-1)^k(\ch\circ h\circ\varphi(x))[T\R^n]\\
&=(-1)^k(\ch\circ\varphi(x))[TN]=(1)^n\big(\psi(\ch(x)\smile\IndexClass(X))\big)[TN]\\
&=(-1)^n(\ch(x)\smile\IndexClass(X))[TX].
\end{align*}

\transnote{上式中的$[T\R^n]$照原书排印。本段的类都是$T\R^k$的基本类，故$[T\R^n]$应读作$[T\R^k]$：$n$是流形$X$的维数，$k$是$X$所嵌入的$\R^k$的维数。}

现在陈述 Atiyah–Singer 指标定理的拓扑版本。

\begin{theorem}[指标定理 A]
若$D$是紧$n$维流形$X$上的椭圆算子，则
\[
\Ind D=(-1)^n\big(\ch(\sigma(D))\smile\IndexClass(X)\big)[TX],
\]
其中$\IndexClass(X)=\Td(TX\otimes_\R\C)\in H^*(X;\Q)$是$X$的指标类。
\end{theorem}

为得到 1963 年指标定理论文中的确切陈述，注意若$X$是定向流形，则对每个$u\in H^n(X;\Z)$有$\psi(u)[TX]=(-1)^{n(n-1)/2}u[X]$，其中$\psi:H^*(X)\to H^*(TX)$是 Thom 同构。因我们对$X$的定向作了选取，必须引入符号$(-1)^{n(n-1)/2}$。于是得到原来的指标定理：

\begin{theorem}[指标定理 B]
若$D$是紧定向$n$维流形$X$上的椭圆算子，则
\[
\Ind D=(-1)^{n(n+1)/2}\big(\psi^{-1}\circ\ch(\sigma(D))\smile\IndexClass(X)\big)[X],
\]
其中$\IndexClass(X)=\Td(TX\otimes_\R\C)\in H^*(X;\Q)$是$X$的指标类。
\end{theorem}

为便于对具体的椭圆算子或复形计算拓扑指标，下面给出 Atiyah–Singer 指标定理的一个特例。在适当情形下，可化简定理陈述中的$\psi^{-1}\circ\ch(\sigma(D))$项。令$\psi:H^*(X)\to H^*(TX)$是 Thom 同构，$i^*:H^*(TX)\to H^*(X)$是零截面$X\to TX$诱导的映射，回想§1：$i^*\psi(x)=x\smile e(X)$，其中$e(X)=e(TX)\in H^n(X;\Z)$是$X$的 Euler 类。于是对一切$y\in H^*(TX)$有$i^*(y)=\psi^{-1}(y)\smile e(X)$，并得到
\[
\psi^{-1}\circ\ch(x)\smile e(X)=i^*(\ch(x))=\ch(i^*(x)).
\]
再设$D^\bullet$是紧流形$X$上的椭圆复形
\[
0\longrightarrow\Gamma E^0\xrightarrow{\ d_0\ }\Gamma E^1\xrightarrow{\ d_1\ }\cdots\xrightarrow{\ d_{n-1}\ }\Gamma E^n\longrightarrow 0,
\]
可见$i^*\sigma(D^\bullet)=\sum_{i=0}^n(-1)^iE^i$，于是
\[
\psi^{-1}\circ\ch(\sigma(D^\bullet))\smile e(X)=\ch\Big(\sum_{i=0}^n(-1)^iE^i\Big).
\]
现在把这个结果并入 Atiyah–Singer 指标定理的陈述。

\transnote{计算链中$(1)^n$一项照原书排印。由上一段的$\ch\circ\varphi(x)=(-1)^{n-k}\psi(\ch(x)\smile\IndexClass(X))$，跨到$[TN]$时应得$(-1)^k\cdot(-1)^{n-k}=(-1)^n$，故此处读作$(-1)^n$才与末行一致。}

% ===== 原书 p.19 =====（§2.3 末尾在 p.19 上半页，已由 c02 处理；本块自 §2.4 起）
\subsection{奇维流形上的指标}

设 $D:\Gamma E\to\Gamma F$ 是紧奇维流形 $X$ 上的椭圆偏微分算子。前面已经指出，椭圆偏微分算子的象征是正齐次映射。用局部坐标写出，对一切实数 $\lambda>0$ 都有 $\sigma(x,\lambda\xi)=\lambda^m\sigma(x,\xi)$，其中 $m$ 是算子的阶。然而象征在局部是多项式，它满足更强的条件：对任意 $\lambda$（包括负值）都有 $\sigma(x,\lambda\xi)=\lambda^m\sigma(x,\xi)$。特别地，$\sigma(x,-\xi)=\pm\sigma(x,\xi)$，其中符号在整个流形上为常数。

于是记 $\alpha:TX\to TX$ 为逐纤维对径映射（丛对合），它由每根纤维上乘以 $-1$ 给出：$\alpha:(x,\xi)\mapsto(x,-\xi)$。考虑诱导映射 $\alpha^*:K(TX)\to K(TX)$。把象征类 $\sigma(D)\in K(TX)$ 用复形
\[
0\to\pi^*E\xrightarrow{\ \sigma(x,\xi)\ }\pi^*F\to 0
\]
表示，则元素 $\alpha^*\sigma(P)$ 由复形
\[
0\to\pi^*E\xrightarrow{\ \sigma(x,-\xi)\ }\pi^*F\to 0
\]
表示。由此得到 $\alpha^*\sigma(P)=\sigma(P)$：若 $\sigma(x,-\xi)=+\sigma(x,\xi)$，这立即成立；若 $\sigma(x,-\xi)=-\sigma(x,\xi)$，则把两个复形沿单位圆旋转半圈，得到它们之间的同伦。

\transnote{「旋转半圈」之所以有效，是因为在复线上 $e^{i\theta}$ 当 $\theta$ 从 $0$ 走到 $\pi$ 时始终可逆，把 $1$ 连续变到 $-1$；若改在实线上做同样的事，途中必然经过 $0$。两个复形因此同伦等价，它们在 $K(TX)$ 中的类相同。}

\begin{theorem}
紧奇维流形上任意椭圆偏微分算子的指标为零。
\end{theorem}

\begin{proof}[证明 1]
回想 Chern 特征给出从有理 K 理论到有理上同调的自然同构，有下面的交换图：
\[
\begin{tikzcd}
K^*(TX)\otimes\Q \arrow[r,"\alpha^*"] \arrow[d,"\ch"{right},"\cong"{left}]
  & K^*(TX)\otimes\Q \arrow[d,"\ch"{right},"\cong"{left}] \\
H^*_c(TX;\Q) \arrow[r,"\alpha^*"']
  & H^*_c(TX;\Q)
\end{tikzcd}
\]
由上同调的 Thom 同构定理，$H^*_c(TX;\Q)$ 是由 Thom 类 $u$ 生成的 $H^*(X;\Q)$ 模；$u$ 是唯一在每个 $TX$ 纤维上限制为选定定向类的上同调类。又因 $X$ 是奇维的，对径映射 $\alpha:TX\to TX$ 反转每根纤维的定向，从而 $\alpha^*u=-u$。

\transnote{图中竖直的 $\ch$ 是 Chern 特征给出的有理同构，两个水平箭头都写作 $\alpha^*$，但分别作用在 K 理论与上同调上。这张图交换，说的正是 $\ch$ 的自然性；证明 1 不过是沿这个方块走一圈。}

% ===== 原书 p.20 =====
由此得到，对每个 $x\in H^*_c(TX;\Q)$ 都有 $\alpha^*x=-x$，因此对应的 K 理论映射也由 $\alpha^*x=-x$ 给出，其中 $x\in K^*(TX)\otimes\Q$。既然 $\alpha^*\sigma(D)=\sigma(D)$，$\sigma(D)$ 在 $K^*(TX)\otimes\Q$ 中的像为零，从而 $\sigma(D)$ 是 $K(TX)$ 中的有限阶元素。拓扑指标是从 $K(TX)$ 到整数的同态，因此 $\Ind D=\tind\sigma(D)=0$。
\end{proof}

\transnote{同态 $\tind:K(TX)\to\Z$ 的像无挠，有限阶元素只能映到 $0$，这是证明 1 的收尾。注意这一步只用到 $\sigma(D)$ 在有理化之后为零，并没有说 $\sigma(D)$ 本身为零。}

\begin{proof}[证明 2]
由指标定理的上同调形式 (A)，有
\[
\Ind D=(-1)^n\big(\ch(\sigma(d))\smile\IndexClass(X)\big)[TX],
\]
其中 $n$ 是 $X$ 的维数。再作用对径映射 $\alpha:TX\to TX$，得到
\[
\begin{aligned}
\Ind D&=(-1)^n\big(\ch(\alpha^*\sigma(D))\smile\IndexClass(X)\big)_{\alpha_*[TX]}\\
&=(-1)^n\big(\ch(\sigma(D))\smile\IndexClass(X)\big)(-1)^n[TX]\\
&=-\Ind D.
\end{aligned}
\]
注意指标类 $\IndexClass(X)=\Td(TX\otimes_\R\C)\in H^*(X,\Q)$ 在 $\alpha^*$ 的作用下不变：复化映射 $\alpha:TX\otimes_\R\C\to TX\otimes_\R\C$ 沿单位圆旋转半圈即同伦于恒等映射。于是 $\Ind D=0$。
\end{proof}

注意，这两个证明的关键事实都是：在任意奇维流形上，对径映射反转 $TX$ 的定向。还要注意，证明 1 说明 $K(TX)$ 的每个元素未必都是某个椭圆偏微分算子的象征。特别地，在圆周上 $K(TS^1)\cong\tilde K(S^2)\cong\Z$，而每个椭圆偏微分算子的象征都是零。

\subsection{de Rham 复形}

现在再来看（复化后的）de Rham 复形 $D^\bullet$：
\[
0\to\Gamma\Lambda^0(TX\otimes\C)\xrightarrow{d}\Gamma\Lambda^1(TX\otimes\C)
\xrightarrow{d}\cdots\xrightarrow{d}\Gamma\Lambda^n(TX\otimes\C)\to 0.
\]
由 de Rham 定理，该复形的欧拉示性数为
\[
\chi(D^\bullet)=\sum_i(-1)^i\dim H^i(D^\bullet)=\sum_i(-1)^i\dim H^i_{dR}(X)=\chi(X).
\]
先考虑 $X$ 为偶数维 $n=2l$ 的情形。为计算拓扑指标，注意由 Euler 类的初等性质有
\[
c_n(TX\otimes\C)=e\big((TX\otimes\C)_\R\big)=(-1)^le(TX\oplus TX)=(-1)^le(X)^2,
\]
其中必须引入符号 $(-1)^l$，因为 $(TX\otimes\C)_\R$ 与 $TX\oplus TX$ 的定向不同。\footnote{设 $\{x_1,\ldots,x_n\}$ 是实向量空间 $V$ 的有序基。取 $\{x_1,ix_1,\ldots,x_n,ix_n\}$ 作为 $(V\otimes\C)_\R$ 的基，取 $\{x_1,\ldots,x_n,x'_1,\ldots,x'_n\}$ 作为 $V\oplus V$ 的基。若 $n=2l$，这两个基之间的置换符号为 $(-1)^l$。}

现在用一个稍作推广的分裂原理，允许（在偶维情形）假设 $TX\otimes\C$ 分解为线丛直和，其形式为
% ===== 原书 p.21 =====
\[
TX\otimes\C=(L_1\oplus\bar L_1)\oplus\cdots\oplus(L_l\oplus\bar L_l).
\]
记 $c_1(L_i)=x_i$，注意 $c_1(\bar L_i)=-x_i$，于是可以算出：
\[
c_n(TX\otimes\C)=\prod_i x_i\smile(-x_i)=(-1)^l\prod_i x_i^2,
\]
\[
\IndexClass(X)=\Td(TX\otimes\C)=\prod_i\frac{x_i}{1-e^{-x_i}}\frac{-x_i}{1-e^{x_i}},
\]
\[
\begin{aligned}
\ch\Big(\sum_{i=0}^n(-1)^i\Lambda^i(TX\otimes\C)\Big)
&=\ch\Big(\prod_i(1-L_i)(1+L_i)\Big)\\
&=\prod_i(1-e^{x_i})(1-e^{-x_i})\\
&=c_n(TX\otimes\C)\smile\Td(TX\otimes\C)^{-1}\\
&=(-1)^le(X)^2\smile\Td(TX\otimes\C)^{-1}.
\end{aligned}
\]
注意下式
\[
\ch\Big(\sum_{i=0}^n(-1)^i\Lambda^i(TX\otimes\C)\Big)\smile e(X)^{-1}
=(-1)^le(X)\smile\Td(TX\otimes\C)^{-1}
\]
在 $e(X)=0$ 时仍良定义，于是可以套用 Atiyah--Singer 指标定理的上同调形式 (C)，得到
\[
\chi(X)=\chi(D^\bullet)=(-1)^{n(n+1)}(-1)^le(X)[X]=e(X)[X].
\]

\transnote{原书此处印作 $(-1)^{n(n+1)}$，而定理 C 中的因子是 $(-1)^{n(n+1)/2}$。取 $n=2l$ 代入可见两者相差 $(-1)^l$，只有取半指数时才有 $\chi(X)=e(X)[X]$。译文照排原文，引用时请以定理 C 为准。}

在奇维情形，Euler 类 $e(X)$ 为零；由上面关于奇维流形上微分算子指标的讨论，$\chi(D^\bullet)=0$ 也成立。于是对任意流形 $X$，de Rham 复形的 Atiyah--Singer 指标定理给出
\[
\xi(X)=e(X)[X].
\]

\subsection{Dolbeault 复形}

现在考察 de Rham 复形的复类比。设 $X$ 是复 $n$ 维流形，记 $T_\C X$ 为 $X$ 的复切丛（即关于 $X$ 上复结构的切丛），有典范同构
\[
TX\otimes\C\cong T_\C X\oplus\overline{T_\C X}.
\]
由此得到分解
\[
\Lambda^i(TX\otimes\C)\cong\sum_{p+q=i}\Lambda^p(\overline{T_\C X})\otimes\Lambda^q(T_\C X)
=\sum_{p+q=i}\Lambda^{p,q}(T_\C X),
\]
其中定义
\[
\Lambda^{p,q}(T_\C X)=\Lambda^p(\overline{T_\C X})\otimes\Lambda^q(T_\C X).
\]
由此可见，$X$ 上复微分形式的空间 $\Gamma\Lambda^i(TX\otimes\C)$ 有分解
\[
\Gamma\Lambda^i(TX\otimes\C)=\sum_{p+q=i}\Gamma\Lambda^{p,q}(T_\C X),
\]
其中 $\Gamma\Lambda^{p,q}(T_\C X)$ 称为 $(p,q)$ 型光滑微分形式的空间。

取 $X$ 上的局部复坐标 $z_1,\ldots,z_n$，其中 $z_j=x_j+iy_j$，$(p,q)$ 型微分形式为
\[
\omega=\sum_{\substack{j_1<\cdots<j_p\\ k_1<\cdots<k_q}}a_{J,K}\,
dz_{j_1}\cdots dz_{j_p}\,d\bar z_{k_1}\cdots d\bar z_{k_q},
\]
这里 $J=(j_1,\ldots,j_p)$ 与 $K=(k_1,\ldots,k_q)$ 是多重指标，$\bar z_j=x_j-iy_j$ 表示复共轭。de Rham 算子 $d:\Gamma\Lambda^i(TX\otimes\C)\to\Gamma\Lambda^{i+1}(TX\otimes\C)$ 于是分解为
% ===== 原书 p.22 =====
$d=\partial+\bar\partial$。这里 $\partial:\Gamma\Lambda^{p,q}(T_\C X)\to\Gamma\Lambda^{p+1,q}(T_\C X)$ 与 $\bar\partial:\Gamma\Lambda^{p,q}(T_\C X)\to\Gamma\Lambda^{p,q+1}(T_\C X)$ 由以下的局部公式定义：
\[
\partial:\sum_{J,K}a_{J,K}dz_Jd\bar z_K\mapsto
\sum_{j,J,K}\frac{\partial}{\partial z_j}a_{J,K}\,dz_jdz_Jd\bar z_K,
\]
\[
\bar\partial:\sum_{J,K}a_{J,K}dz_Jd\bar z_K\mapsto
\sum_{k,J,K}\frac{\partial}{\partial\bar z_k}a_{J,K}\,d\bar z_kdz_Jd\bar z_K,
\]
其中
\[
\frac{\partial}{\partial z_j}=\frac12\Big(\frac{\partial}{\partial x_j}-i\frac{\partial}{\partial y_j}\Big),
\qquad
\frac{\partial}{\partial\bar z_j}=\frac12\Big(\frac{\partial}{\partial x_j}+i\frac{\partial}{\partial y_j}\Big).
\]
把 de Rham 复形限制到 $(0,q)$ 型微分形式，得到 Dolbeault 复形，它是如下椭圆复形：
\[
0\to\Gamma\Lambda^{0,0}(T_\C X)\xrightarrow{\bar\partial}\Gamma\Lambda^{0,1}(T_\C X)
\xrightarrow{\bar\partial}\cdots\xrightarrow{\bar\partial}\Gamma\Lambda^{0,n}(T_\C X)\to 0.
\]

为计算 Dolbeault 复形的指标，注意 $X$ 是复流形，Euler 类为 $e(X)=e(TX)=c_n(T_\C X)$。因此可以套用指标定理的形式 (C)。按定义有 $\Lambda^{0,q}(T_\C X)=\Lambda^q(T_\C X)$。

\transnote{这一步是定义层面的恒等：$\Lambda^{0,q}(T_\C X)=\Lambda^0(\overline{T_\C X})\otimes\Lambda^q(T_\C X)$，而 $\Lambda^0$ 是平凡线丛，故只剩 $\Lambda^q(T_\C X)$。不要把 $\Lambda^{0,q}$ 读成外代数里某个新的直和项。}

用分裂原理，不妨设 $T_\C X=L_1\oplus\cdots\oplus L_n$，$x_i=c_1(L_i)$，算出：
\[
\ch\Big(\sum_{i=0}^n\Lambda^i(T_\C X)\Big)=\ch\Big(\prod_i(1-L_i)\Big)=\prod_i(1-e^{x_i}),
\]
\[
e(X)=c_n(T_\C X)=\prod_i x_i,
\]
\[
\IndexClass(X)=\Td(TX\otimes\C)=\Td(T_\C X)\smile\Td(\overline{T_\C X})
=\Td(T_\C X)\smile\prod_i\frac{-x_i}{1-e^{x_i}}.
\]
把这些表达式代入指标定理，得到
\[
\begin{aligned}
\chi(\bar\partial)&=(-1)^{2n(2n+1)/2}
\Big(\ch\Big(\sum_{i=0}^n(-1)^i\Lambda^i(T_\C X)\Big)\smile e(X)^{-1}\smile\Td(TX\otimes\C)\Big)[X]\\
&=(-1)^{n(2n+1)}
\Big(\prod_i\frac{1-e^{x_i}}{x_i}\smile\Td(T_\C X)\smile\prod_i\frac{-x_i}{1-e^{x_i}}\Big)[X]\\
&=(-1)^n\big((-1)^n\Td(T_\C X)\big)[X]=\Td(T_\C X)[X].
\end{aligned}
\]
推广这一构造，可以得到著名的 Hirzebruch--Riemann--Roch 定理。完整的讨论见 [11, \S4]。

\begin{example}
把球面 $S^2$ 看作复 $1$ 维流形（即 Riemann 球），Dolbeault 复形约化为 $S^2$ 上的 Dolbeault 算子
\[
\bar\partial:\Gamma\Lambda^0(T_\C S^2)\to\Gamma\Lambda^1(T_\C S^2),
\]
它由 $\bar\partial f=(\partial/\partial\bar z)\,d\bar z$ 给出。由上面的讨论有
\[
\Ind\bar\partial=\Td(T_\C S^2)[X].
\]
现在回想 Todd 类由
\[
\Td(E)=1+\frac12c_1(E)+(\text{高阶项})
\]
给出，
% ===== 原书 p.23 =====
代入后得到
\[
\begin{aligned}
\Ind\bar\partial&=\Td(T_\C S^2)[X]=\big((1+\tfrac12c_1(T_\C S^2)\big)[X]\\
&=\tfrac12\big(c_1(T_\C S^2)\big)[X]=\tfrac12e(X)[X],
\end{aligned}
\]
而上一步刚证明 $e(X)[X]=\chi(X)$，于是
\[
\Ind\bar\partial=\frac12\chi(S^2)=1.
\]
更一般地，设 $X$ 是亏格 $g$ 的 Riemann 面，则 $\chi(X)=2-2g$，从而
\[
\Ind\bar\partial_X=1-g.
\]
\end{example}

% ===== 原书 p.23 =====
\section{指标定理的证明}

\subsection{拓扑指标的公理}

由拓扑指标的构造立即看出，它满足下面两条初等公理：

\noindent\textbf{公理 A}\quad 若 $X$ 是一点，则 $\tind_{\pt}:\Z\to\Z$ 是恒等映射。

\noindent\textbf{公理 B}\quad $\tind$ 与同态 $i_!$ 交换。

事实上，后面将看到，这两条公理唯一刻画了拓扑指标。

\begin{definition}
对每个紧流形 $X$ 给定一族同态 $\operatorname{ind}_X:K(TX)\to\Z$。若它关于微分同胚是函子式的，就称为一个指标函数（index function）。换言之，若 $f:X\to Y$ 是微分同胚，则对一切 $y\in K(TY)$ 都有 $\operatorname{ind}_Xf^*y=\operatorname{ind}_Yy$，其中 $f^*:K(TY)\to K(TX)$ 是由 $f$ 到切丛的延拓诱导的映射。
\end{definition}

\begin{proposition}
若 $\operatorname{ind}$ 是满足上述公理 A 与公理 B 的指标函数，则 $\operatorname{ind}=\tind$。
\end{proposition}

\begin{proof}
对任意紧流形 $X$，取嵌入 $i:X\to\R^m$，并令 $j:\pt\to\R^m$ 为原点的包含映射。由公理 A 可知下图交换：
\[
\begin{tikzcd}
K(TX) \arrow[r,"i_!"] \arrow[dr,"\operatorname{ind}"']
  & K(T\R^m) \arrow[r,"j_!^{-1}"] & K(T\pt) \arrow[dl,"\operatorname{ind}"] \\
& \Z & {}
\end{tikzcd}
\]
于是得到
\[
\operatorname{ind}_X=\operatorname{ind}_{\pt}\circ j_!^{-1}\circ i_!
=j_!^{-1}\circ i_!=\tind_X.
\]
\end{proof}

\transnote{$j_!^{-1}$ 之所以存在，是因为 $T\R^m$ 可缩，$K(T\R^m)$ 由原点附近生成，$j_!$ 从 $K(T\pt)$ 到 $K(T\R^m)$ 恰好是同构。公理 A 把单点上的值钉死，公理 B 让推前可以穿过 $\tind$；两条合起来，任何 $X$ 上的 $\tind$ 都由嵌入与单点决定。}

Atiyah--Singer 指标定理的证明方向现在清楚了。我们想构造一个解析指标，在象征层面上由 $\aind=\Ind\circ\sigma^{-1}$ 给出：把每个元素 $x\in K(TX)$ 映到某个象征为 $x$ 的椭圆算子的指标。于是指标定理归结为证明拓扑指标与解析指标相等，为此只需对解析指标验证公理 A 与公理 B。此时公理 B 告诉我们：给定 $X$ 上的椭圆算子 $D$ 与包含映射 $i:X\to Y$，可以（在象征层面）构造
% ===== 原书 p.24 =====
$Y$ 上的椭圆算子 $i_!D$，它与 $D$ 有相同的指标。于是，既然任何流形都能嵌入 $\R^m$，指标问题就化归为 $m$ 球面 $S^m=(\R^m)^+$ 上椭圆算子的问题，而后者容易解决得多。可惜，正如 §2 讨论奇维流形时看到的，$K(TX)$ 的每个元素未必是某个椭圆偏微分算子的象征。为弥补这一点，需要把讨论推广到伪微分算子类。

\subsection{伪微分算子}

设 $f$ 是 $\R^n$ 上具有紧支撑的实值函数。把 $f$ 的 Fourier 变换定义为函数 $\hat f$：
\[
\hat f(\xi)=\frac{1}{(2\pi)^n}\int_{\R^n}f(x)e^{-i\langle x,\xi\rangle}dx.
\]
Fourier 反演公式把 $f$ 用 $\hat f$ 写成积分：
\[
f(x)=\int_{\R^n}\hat f(\xi)e^{i\langle x,\xi\rangle}.
\]
取 $f$ 的偏导数 $D_j=-i\partial/\partial x_j$，得到
\[
-i\frac{\partial}{\partial x_j}f(x)
=\int_{\R^n}-i\frac{\partial}{\partial x_j}\hat f(\xi)e^{i\langle x,\xi\rangle}d\xi
=\int_{\R^n}\xi_j\hat f(\xi)e^{i\langle x,\xi\rangle}d\xi.
\]
于是，Fourier 变换「把微分变成乘法」。特别地，若 $P$ 是 $\R^n$ 上的偏微分算子，则有
\[
Pf(x)=\int_{\R^n}p(x,\xi)\hat f(\xi)e^{i\langle x,\xi\rangle}d\xi,
\]
其中 $p(x,\xi)$ 是 $\xi=(\xi_1,\ldots,\xi_n)$ 的多项式，其系数是 $x$ 的光滑实值函数。

为看清 $P$ 与 $p(x,\xi)$ 的联系，可以写 $P=p(x,D)$，其中 $D=(-i\partial/\partial x_1,\ldots,-i\partial/\partial x_n)$。注意可以通过取交换子从算子 $P$「提取」多项式 $p(x,\xi)$：
\[
p(x,\xi)=e^{-i\langle x,\xi\rangle}Pe^{i\langle x,\xi\rangle}.
\]
跃迁到伪微分算子，发生在允许 $p(x,\xi)$ 不必是多项式的时候。此时只要求 $p(x,\xi)$ 满足适当的增长条件，并且能够定义象征 $\sigma(x,\xi)$。

\begin{definition}
从 $\R^n$ 上具有紧支撑的光滑函数到 $\R^n$ 上光滑函数的线性算子 $P$，若由 $P=p(x,D)$ 给出，就称为 $m$ 阶\textbf{伪微分算子}（pseudo-differential operator），其中 $p(x,\xi)$ 是满足增长条件
\[
\big|D_x^\beta D_\xi^\alpha p(x,\xi)\big|\le C_{\alpha,\beta}\big(1+|\xi|\big)^{m-|\alpha|}
\]
的光滑函数。\footnote{这些增长条件允许在积分号下求导。}
\end{definition}

% ===== 原书 p.25 =====
其中对每个多重指标 $\alpha$、$\beta$ 有正常数 $C_{\alpha,\beta}$。这里记 $|\alpha|=\alpha_1+\cdots+\alpha_n$，$D_\xi^\alpha=(\partial/\partial\xi_1)^{\alpha_1}\cdots\allowbreak(\partial/\partial\xi_n)^{\alpha_n}$。还要求对一切 $x\in\R^n$ 与 $\xi\neq 0$，极限
\[
\sigma_m(x,\xi)=\lim_{\lambda\to\infty}\frac{p(x,\lambda\xi)}{\lambda^m}
\]
存在；$\sigma_m(x,\xi)$ 称为 $P$ 的 $m$ 阶\textbf{象征}（symbol）。

若 $P$ 是偏微分算子，则对应的多项式 $p(x,\xi)$ 满足增长条件，$\sigma_m(x,\xi)$ 等于去掉 $p(x,\xi)$ 中除最高阶项外所有项后定义的象征。还要注意，阶 $m$ 不必是整数。可以取 $m=\frac12$，例如 $P=p(x,D)$，其中 $p(x,\xi)=a_1(x)\sqrt{\xi_1}+\cdots+a_n(x)\sqrt{\xi_n}$；也可以考虑 $m\le 0$ 的算子。无论如何，$P$ 的象征 $\sigma_m(x,\xi)$ 都有性质
\[
\sigma_m(x,\lambda\xi)=\lambda^m\sigma_m(x,\lambda\xi)
\quad\text{对一切实数 }\lambda>0,
\]
于是称 $\sigma_m(x,\xi)$ 为 $m$ 次\textbf{正齐次的}（positively homogeneous）。

\transnote{原文此处印作 $\sigma_m(x,\lambda\xi)=\lambda^m\sigma_m(x,\lambda\xi)$，右端应为 $\sigma_m(x,\xi)$。极限定义 $\sigma_m(x,\xi)=\lim_{\lambda\to\infty}p(x,\lambda\xi)/\lambda^m$ 本身就迫使 $\sigma_m$ 关于 $\xi$ 正齐次，读者按此理解即可。译文照排原文。}

这个定义带来一个问题：函数 $e^{i\langle x,\xi\rangle}$ 没有紧支撑，因此上面那个交换子无法取。为补救，把定义局部化：只要求对每个紧支撑函数 $f$，存在满足上述增长条件的光滑函数 $p_f(x,\xi)$，使得对任意函数 $u$（支撑任意）都有
\[
P(f\cdot u)=p_f(x,D)u.
\]
此时可以取交换子
\[
p_f(x,\xi)=e^{-\langle x,\xi\rangle}P\big(f\cdot e^{i\langle x,\xi\rangle}\big).
\]
于是可以这样定义 $P$ 的象征：$\sigma(x,\xi)=\sigma_f(x,\xi)$，其中 $f$ 是紧支撑函数，且在 $x$ 的某个邻域内等于 $1$。结果表明这个定义与 $f$ 的选取无关。

我们的定义容易推广到向量值函数 $f:\R^m\to\R^n$ 上的伪微分算子；此时多项式 $p_f(x,\xi)$ 与象征 $\sigma(x,\xi)$ 是 $x$ 与 $\xi$ 的矩阵值函数。若伪微分算子 $P$ 的象征 $\sigma(x,\xi)$ 对一切 $x$ 与 $\xi\neq0$ 都可逆，就称 $P$ 是\textbf{椭圆的}（elliptic）。现在可以把伪微分算子的概念推广到流形。

\begin{definition}
设 $E$ 与 $F$ 是流形 $X$ 上的复向量丛。若线性算子 $P:\Gamma_cE\to\Gamma F$ 满足：存在 $X$ 的坐标卡覆盖 $\{U_u\}$，在其上 $E$、$F$、$TX$ 都平凡，且把 $P$ 限制到 $U_i$ 中紧支撑函数得到的算子 $P_i$ 是 $\R^n$ 上的 $m$ 阶伪微分算子（$n=\dim X$），则称 $P$ 是 $X$ 上 $m$ 阶伪微分算子。\footnote{这里用记号 $\Gamma_c$ 表示紧支撑光滑截面的空间。}
\end{definition}

于是可以定义整体象征 $\sigma(P):\pi^*E\to\pi^*F$。若伪微分算子 $P$ 的象征 $\sigma(P)$ 在零截面之外可逆，就称 $P$ 是椭圆的。如前，若 $X$ 紧，象征给出元素 $\sigma(P)\in K(TX)$。事实上，$K(TX)$ 的任一元素都可以这样得到。在证明这一点之前，先算一个例子。

% ===== 原书 p.26 =====
\begin{example}
§2 已证明，奇维流形上任意椭圆偏微分算子的指标为零。特别地，圆周上这样的算子指标为零。现在构造圆周上一个指标为 $-1$ 的椭圆伪微分算子。考虑圆周 $S^1=\R/2\pi\Z$ 上复值函数的线性算子 $P$，它按 Fourier 级数展开给出：
\[
P(e^{ik\theta})=
\begin{cases}
e^{ik\theta} & k\ge 0,\\
0 & k<0.
\end{cases}
\]
于是 $P$ 实际上杀掉了 Fourier 级数展开中 $k<0$ 的部分。给定任一在区间 $0<\theta<2\pi$ 中紧支撑的复值函数 $f$，得到
\[
\begin{aligned}
p_f(\theta,\xi)&=e^{-i\theta\xi}P\big(f\cdot e^{i\theta\xi}\big)
=e^{-i\theta\xi}P\Big(\sum_{-\infty}^{\infty}\hat f(k)e^{ik\theta}\cdot e^{i\theta\xi}\Big)\\
&=e^{-i\theta\xi}P\Big(\sum_{-\infty}^{\infty}\hat f(k-\xi)e^{ik\theta}\Big)
=\sigma_0^{\infty}\hat f(k-\xi)e^{i\theta(k-\xi)},
\end{aligned}
\]

\transnote{此式最右端原书印作 $\sigma_0^{\infty}\hat f(k-\xi)e^{i\theta(k-\xi)}$，其中 $\sigma_0^{\infty}$ 应是求和号 $\sum_{k=0}^{\infty}$。否则 $k$ 没有约束，这一步也与前一行的 $\sum_{-\infty}^{\infty}$ 不衔接。译文照排原文，理解时按求和处理。}

取极限即见
\[
p_f(\theta,\xi)\to
\begin{cases}
f(\theta) & \text{当 }\xi\to+\infty,\\
0 & \text{当 }\xi\to-\infty.
\end{cases}
\]
因此 $p_f(\theta,\xi)$ 满足零阶伪微分算子的增长条件，整体零阶象征为
\[
\sigma_p(\theta,\xi)=
\begin{cases}
1 & \xi>0,\\
0 & \xi<0.
\end{cases}
\]
这个算子不是椭圆的，但算子 $A=e^{i\theta}P+(1-P)$ 也是零阶伪微分算子；由于它的象征是
\[
\sigma_A(\theta,\xi)=
\begin{cases}
e^{i\theta} & \xi>0,\\
1 & \xi<0,
\end{cases}
\]
可见 $A$ 是椭圆的。此外，按 Fourier 级数展开有
\[
A(e^{ik\theta})=
\begin{cases}
e^{i(k+1)\theta} & k\ge 0,\\
e^{ik\theta} & k<0.
\end{cases}
\]
于是 $A$ 把 Fourier 级数中 $k\ge 0$ 的部分挪动一项；又因 $\Ker A$ 为空，而 $\Coker A$ 可以等同于常数函数空间，所以 $\Ind A=-1$。这个例子在 Bott 周期性定理证明中的推广见附录。
\end{example}

\transnote{算子 $A$ 的象征在 $\xi>0$ 处取 $e^{i\theta}$、在 $\xi<0$ 处取 $1$，两个分支都非零，故 $A$ 椭圆。微分算子的象征在 $\xi$ 与 $-\xi$ 处只差一个 $\pm$ 号，做不出两个分支各取一个可逆值的效果；奇维流形上的指标定理因此必须放下微分算子，改用伪微分算子。}

\subsection{解析指标的构造}

\begin{claim}
给定 $a\in K(TX)$，其中 $X$ 是紧流形，则存在 $X$ 上任一给定阶 $m$ 的椭圆伪微分算子 $P$，使得 $\sigma(P)=a$。
\end{claim}

\begin{proof}
已知 $a$ 可以由长度为一的复形
\[
0\to\pi^*E\xrightarrow{\alpha_m}\pi^*F\to 0,
\]
表示，
% ===== 原书 p.27 =====
其中 $\alpha_m$ 是 $m$ 次齐次的。取使 $E$、$F$、$TX$ 平凡的局部坐标，令 $p(x,\xi)=\varphi(\xi)\alpha_m(x,\xi)$，其中 $\varphi(\xi)$ 是光滑函数，在原点的邻域内为零，其余各处为 $1$。这个修正项是必要的：当 $m\le 0$ 时，$\alpha_m$ 可能在零截面处不连续。于是直接看出，$P=p(x,D)$ 是 $X$ 上 $m$ 阶伪微分算子，其象征恰为 $\alpha_m$。从而 $P$ 是椭圆算子，且 $\sigma(P)=a$。
\end{proof}

既然已说明象征映射是满射，我们还想说明象征相同的两个椭圆算子指标相同（这样就能说象征映射是「伪单射」）。在证明这一点之前，必须先确认椭圆伪微分算子的指标确有定义。所幸，若 $P$ 是椭圆伪微分算子，则由分析可知 $\Ker P$ 与 $\Coker P$ 都有限维\footnote{这远不是显然的结论。关于 Sobolev 空间的讨论见 [10, \S5]。若仍不放心，可以与做分析的人讨论这个问题。}，于是可以把指标定义为
\[
\Ind P=\dim\Ker P-\dim\Coker P.
\]
还知道 $m$ 阶椭圆伪微分算子的指标只依赖象征在同阶象征空间中的同伦类。现在可以证明：

\begin{claim}
若 $P$ 与 $P'$ 是紧流形 $X$ 上两个椭圆伪微分算子，且 $\sigma(P)=\sigma(P')\in K(TX)$，则 $\Ind P=\Ind P'$。
\end{claim}

\begin{proof}
设 $\sigma(P)$ 与 $\sigma(P')$ 分别由复形
\[
0\to\pi^*E\xrightarrow{\sigma_m}\pi^*F\to 0,
\qquad
0\to\pi^*E'\xrightarrow{\sigma'_n}\pi^*F'\to 0
\]
给出，其中 $\sigma_m$ 与 $\sigma'_n$ 分别是 $P$ 与 $P'$ 的象征。先设 $P$ 与 $P'$ 同为 $m$ 阶。既然 $\sigma(P)=\sigma(P')$，可见这两个复形或者同伦等价，或者相差一个支撑为空的复形。若同伦，则 $\sigma_m$ 与 $\sigma'_m$ 通过 $m$ 次齐次同态同伦，由此得到 $\Ind P=\Ind P'$。若二者相差形如
\[
0\to G\xrightarrow{\beta}G\to 0
\]
的复形，其中 $\beta$ 是同构；由于 $\beta$ 也是 $m$ 次齐次的，当 $m\neq0$ 时必有 $\dim G=0$。当 $m=0$ 时，$\beta$ 关于 $\xi$ 为常数，于是 $P$ 与 $P'$ 相差算子 $Qu=\beta(x,0)u$，且 $\Ind Q=0$。同样得到 $\Ind P=\Ind P'$。

现设 $P$ 是 $m$ 阶而 $P'$ 是 $n$ 阶，$m\neq n$。由上面的论证，不妨设 $E=E'$、$F=F'$，且 $\sigma_m=\sigma'_n$ 在 $S(X)$ 上成立。构造映射 $\rho=\sigma'_n\sigma_m^{-1}$，可见 $\rho$ 是 $n-m$ 次齐次的，且在 $S(X)$ 上是恒等映射。于是 $\rho$ 自伴；由此 $P$ 与 $P'$ 相差一个相伴的自伴算子 $R$。由于 $R$ 自伴，有 $\Ind R=\dim\Ker R-\dim\Ker R^*=0$，从而 $\Ind P=\Ind P'$。
\end{proof}

\transnote{$R$ 自伴时 $\Coker R\cong\Ker R^*=\Ker R$，故 $\Ind R=\dim\Ker R-\dim\Ker R^*=0$。这一步默许 $\Ker R$ 与 $\Ker R^*$ 都有限维，依据是脚注里引用的分析结论。}

\begin{definition}
\textbf{解析指标}（analytical index）是映射 $\aind:K(TX)\to\Z$，由
\[
\aind:x\mapsto\Ind P,\qquad\text{其中 }\sigma(P)=x
\]
给出。
\end{definition}

% ===== 原书 p.28 =====
\subsection{公理的验证}

由构造容易看出，a-ind 关于微分同胚是函子的，因此 a-ind 是指标函数。下面要说明解析指标满足两条公理：

\noindent\textbf{公理 A}\quad 若 $X$ 是一点，则 $\aind\pt:\Z\to\Z$ 是恒等映射。

\noindent\textbf{公理 B}\quad a-ind 与同态 $i_!$ 交换。

这两条公理唯一刻画拓扑指标 t-ind。于是只要它们成立，就得到 a-ind = t-ind，Atiyah–Singer 指标定理随即成立。

验证公理 A。若 $X$ 是一点，则 $X$ 上的椭圆算子只是一个线性变换 $P:V\to W$，其中 $V$、$W$ 是复向量空间，并且
\[
\sigma(P)=\dim E-\dim F\in K(TX)\cong\Z .
\]
此外注意
\[
\dim\Ker P+\dim\operatorname{Im}P=\dim E,\qquad
\dim F-\dim\operatorname{Im}P=\dim\Coker P .
\]
于是
\[
\Ind P=\dim\Ker P-\dim\Coker P=\sigma(P),
\]
公理 A 立即成立。

\transnote{这两条公理把指标定理拆成归一点与切除两步。要点在唯一性：任何满足两条公理的指标函数都等于 t-ind，t-ind 又按构造满足它们。所以证明 a-ind = t-ind 不必直接计算 a-ind，只需验证公理。}

验证公理 B 要麻烦一些。回忆紧子流形 $X\subset Y$ 的包含映射 $i:X\to Y$，我们曾把 $i_!$ 定义为合成
\[
i_!:K(TX)\xrightarrow{\ \varphi\ }K(TN)\xrightarrow{\ h_*\ }K(TY),
\]
其中 $N$ 是 $X$ 在 $Y$ 中的管状邻域，$\varphi$ 是 Thom 同构，$h_*$ 是由开包含 $h:TN\to TY$ 诱导的自然延拓同态。要验证 a-ind 与 $i_!$ 交换，只须证明它与 $\varphi$ 及 $h_*$ 都交换。先看自然延拓同态 $h_*$。

\begin{proposition}[切除]
若 $h:U\to X$ 是开集 $U$ 到 $X$ 的包含映射，则
\[
\aind\circ h_*(x)=h_*\circ\aind(x),
\]
其中 $h_*:K(TU)\to K(TX)$ 是自然延拓同态。
\end{proposition}

\begin{proof}
$K(TU)$ 中任一元素 $a$ 可由 $TU$ 上的齐次复形
\[
0\longrightarrow\pi^*E\xrightarrow{\ \alpha\ }\pi^*F\longrightarrow 0
\]
表示。$U$ 未必紧，故须取 $\alpha$ 为零次齐次的；还要求 $\alpha(x,\xi)=\mathrm{Id}$ 对所有位于某个紧集 $C\subset U$ 之外的 $x$ 成立（非紧情形的讨论见 [10, pp. 493–3]）。把 $E$、$F$ 平凡地延拓为 $X$ 上的丛 $E'$、$F'$，再考虑 $X$ 上的复形
\[
0\longrightarrow\pi^*E'\xrightarrow{\ \alpha\ }\pi^*F'\longrightarrow 0,
\]
% ===== 原书 p.29 =====
它表示元素 $h_*a\in K(TX)$，其中
\[
\alpha'=
\begin{cases}
\alpha, & \text{在 }TU\text{ 上},\\
\mathrm{Id}, & \text{在 }TX-TU\text{ 上}.
\end{cases}
\]
设 $P$、$P'$ 分别是 $U$ 与 $X$ 上对应的算子。由构造有 $\sigma(P')=h_*\sigma(P)$。若 $f$ 的支撑在 $U$ 内且 $Pf=0$，则 $f$ 的支撑必在 $U$ 内，从而 $P'f=0$。于是 $\Ker P=\Ker P'$。对伴算子 $P^*$ 与 $P'^*$ 同理，所以 $\Ind P=\Ind P'$。
\end{proof}

\transnote{原文取零次齐次并要求 $\alpha$ 在紧集外为恒等，是紧支撑 K 理论的约定：只有象征在无穷远处趋于恒等，才能把 $E$、$F$ 平凡延拓到 $X$ 而不改变 $K(TX)$ 中的类。若省掉紧支撑条件，延拓后的复形未必还表示 $h_*a$。}

剩下要证的是：解析指标与 Thom 同构 $\varphi:K(X)\to K(V)$ 交换，其中 $V$ 是 $X$ 上的复向量丛。回忆 $\varphi(x)=\pi^*x\cdot\lambda_V$，其中 $\pi:V\to X$ 是丛投影，$\lambda_V=\varphi(1)$ 是由外复形 $\Lambda(V)$ 给出的 Thom 类。要证的是
\[
\aind(\pi^*x\cdot\lambda_V)=\aind(x)\cdot\aind(\lambda_n)=\aind(x).
\]
为此须先建立 a-ind 的乘积公式，再在适当条件下验证 $\aind(\lambda_n)=1$。先从乘积公式入手。若丛可写成平凡丛 $V=X\times\R^n$，问题会大大简化。对一般向量丛，须改而考虑「扭」乘积。

设有主纤维丛 $P\to X$，其相伴的紧 Lie 群为 $H$：$H$ 从右自由作用在 $P$ 上，$X=P/H$。再设 $F$ 是局部紧空间，$H$ 从左作用在 $F$ 上，则得到 $X$ 上的纤维丛 $Y$：
\[
Y=P\times_H F=(P\times F)\big/\big((p,f)\sim(ph^{-1},hf)\ \ \forall h\in H\big).
\]
$H$ 在 $F$ 上的作用诱导 $TF$ 上的 $H$ 作用，因此可构造 $Y$ 上的向量丛 $P\times_H TF$。取度量后得到分解
\[
TY=(P\times_H TF)\oplus\pi^*TX,
\]
从而有乘积映射 $K(TX)\otimes K(P\times_H TF)\to K(TY)$。考虑同态
\[
K_H(TF)\to K_H(P\times TF)\cong K(P\times_H TF),
\]
其中 $K_H(\cdot)$ 是 §4 讨论的等变 K 理论函子，于是得到乘法映射
\[
K(TX)\otimes K_H(TF)\to K(TY).
\]
我们关心 $H=O(n)$、$F=\R^n$ 的情形。此时任一向量丛 $Y\to X$ 都可写成 $Y=P\times_{O(n)}\R^n$，其中 $P\to X$ 是适当的主纤维丛；于是得到乘法映射
\[
K(TX)\otimes K_{O(n)}(T\R^n)\to K(TY).
\]
下面给出解析指标的乘积公式。

\begin{proposition}[乘性公理]
若 $a\in K(TX)$，$b\in K_H(TF)$，则
\[
\aind^Y a\cdot b=\aind^H a\cdot\aind_H^F b,
\]
只要 $\aind_H^F b\in\Z$ 是平凡表示 $1\in R(H)$ 的整数倍。这里对 $\sigma(P)=b$ 定义 $\aind_H b=\Ind_H P=[\Ker P]-[\Coker P]\in R(H)$（见 §4）。
\end{proposition}

\transnote{等号左边省略了括号：$\aind^Y a\cdot b$ 指 $\aind^Y(a\cdot b)$。$a$ 是 $K(TX)$ 中的类，本身不在 $K(TY)$ 中，只有乘积 $a\cdot b$ 才有指标。}

% ===== 原书 p.30 =====
\begin{proof}
元素 $a$ 可由 $TX$ 上的 1 次齐次复形
\[
0\longrightarrow\pi^*E^0\xrightarrow{\ \alpha\ }\pi^*E^1\longrightarrow 0
\]
表示。构造 $X$ 上以 $\alpha$ 为象征的伪微分椭圆算子 $A$。取从属于平凡化 $Y$ 的开覆盖 $\{U_i\}$ 的单位分解，可把 $A$ 提升为 $Y$ 上的整体椭圆算子 $\tilde A$。\footnote{若 $P:\Gamma E\to\Gamma F$ 是 $X$ 上的算子，$G$ 是 $Y$ 上的向量丛，则 $P$ 到 $X\times Y$ 的提升（lift）$\tilde P:\Gamma(E\BT G)\to\Gamma(F\BT G)$ 定义为 $\tilde P(u(x)\otimes v(y))=P(u(x))\otimes v(y)$。}同理，$b$ 可由 $TF$ 上的 1 次齐次复形
\[
0\longrightarrow\pi^*G^0\xrightarrow{\ \beta\ }\pi^*G^1\longrightarrow 0
\]
表示。

设 $B$ 是 $F$ 上以 $\beta$ 为象征的椭圆伪微分算子，并要求 $B$ 与 $H$ 在 $G^0$、$G^1$ 上的作用交换。于是 $B$ 可提升为 $Y$ 上的整体椭圆算子 $\tilde B$。

现在构造 $Y$ 上的算子
\[
D=\begin{pmatrix}\tilde A & -\tilde B^*\\ \tilde B & \tilde A^*\end{pmatrix},
\qquad\text{其象征为}\qquad
\sigma(D)=\begin{pmatrix}\alpha\BT 1 & -1\BT\beta^*\\ 1\BT\beta & \alpha^*\BT 1\end{pmatrix}.
\]
回忆两个长度为 1 的复形之积的表达式，可知 $\sigma(D)$ 对应 $K(TY)$ 中的元素 $ab$，它由复形
\[
0\longrightarrow E^0\BT G^0\oplus E^1\BT G^1\xrightarrow{\ \sigma(D)\ }E^1\BT G^0\oplus E^0\BT G^1\longrightarrow 0
\]
表示。因此只须计算 $D$ 的指标。注意 $\tilde A$ 与 $\tilde B$ 交换，考察两个对角化算子
\[
D^*D=\begin{pmatrix}\tilde A^*\tilde A+\tilde B^*\tilde B & 0\\ 0 & \tilde A\tilde A^*+\tilde B\tilde B^*\end{pmatrix}
=\begin{pmatrix}P_0 & 0\\ 0 & Q_0\end{pmatrix},
\]
\[
DD^*=\begin{pmatrix}\tilde A\tilde A^*+\tilde B^*\tilde B & 0\\ 0 & \tilde A^*\tilde A+\tilde B\tilde B^*\end{pmatrix}
=\begin{pmatrix}P_1 & 0\\ 0 & Q_1\end{pmatrix},
\]
其中
\[
\Ker D=\Ker D^*D=\Ker P_0\oplus\Ker Q_0,
\]
\[
\Coker D=\Ker D^*=\Ker DD^*=\Ker P_1\oplus\Ker Q_1.
\]
于是
\[
\Ind D=(\dim\Ker P_0-\dim\Ker P_1)+(\dim\Ker Q_0-\dim\Ker Q_1).
\]
考察算子 $P_0=\tilde A^*\tilde A+\tilde B^*\tilde B$，注意
\[
\ip{P_0 u}{u}=\ip{\tilde A u}{\tilde A u}+\ip{\tilde B u}{\tilde B u},
\]
所以 $\Ker P_0=\Ker\tilde A\cap\Ker\tilde B$。$\tilde B$ 把 $B$ 延拓到 $Y$ 的各纤维，因此 $\Ker\tilde B$ 是 $X$ 上向量丛 $K_B=P\times_H\Ker B$ 的光滑截面空间。$\tilde A$ 诱导 $K_B$ 的截面上的算子 $C$，其象征 $\sigma(C)=\alpha\otimes\mathrm{Id}(K_B)$，于是 $\sigma(C)=a[K_B]\in K(TX)$，

% ===== 原书 p.31 =====
其中 $[K_B]$ 是 $K_B$ 在 $K(X)$ 中的类。把 $\tilde A$ 换成 $\tilde A^*$，$P_1$ 与 $C^*$ 有类似结论，于是
\[
\begin{aligned}
\dim\Ker P_0-\dim\Ker P_1
&=\dim\Ker C-\dim\Ker C^*\\
&=\Ind C=\aind^X(a[K_B]).
\end{aligned}
\]
同样，取 $L_B=P\times_H\Coker B$，得到
\[
\dim\Ker Q_0-\dim\Ker Q_1=\aind^X(a[L_B]).
\]
合并两式，得
\[
\Ind D=\aind^X\big(a\cdot([K_B]-[L_B])\big).
\]
若还有 $\aind_H^F b=[\Ker B]-[\Coker B]\in R(H)$ 是整数（即平凡表示 $\C$ 的整数倍），则 $[K_B]-[L_B]=(\aind_H^F b)\cdot 1$ 是平凡丛 $1$ 的整数倍。又因 a-ind$^X$ 是同态，即得所求的乘积公式
\[
\Ind^Y D=(\aind^X a)\cdot(\aind_H^F b).
\]
\end{proof}

\begin{proposition}[归一化公理]
若 $j:\pt\to\R^n$ 是原点的包含映射，于是有诱导同态 $j_!:R(O(n))\to K_{O(n)}(T\R^n)$，则
\[
\aind_{O(n)}j_!(1)=\aind_{O(n)}\lambda_n=1,
\]
其中 Bott 类 $\lambda_n=\lambda_{\C^n}$ 是 $\C^n$ 的 K 理论 Thom 类。
\end{proposition}

\begin{proof}
先看 $n=2$。回忆 de Rham 象征 $\rho_2=\rho_{S^2}$ 由 $2$-球面 $S^2$ 的 de Rham 复形复化得到。涉及的丛可定向，故结构群可取为 $SO(2)$。考察 $SO(2)$ 在 $S^2=\R^2\cup\infty$ 上的作用，可知 $\rho_2\in K_{SO(2)}(TS^2)$；经过适当形变得到（见 [10, §3.2]）
\[
\rho_2=h(\lambda_2)+f^*\circ h(\lambda_2),
\]
其中 $h$ 是延拓同态 $K(T\R^2)\to K(TS^2)$，$f:S^2\to S^2$ 是交换 $0$ 与 $\infty$ 的反射。a-ind 是函子的，且与自然延拓同态交换，故 $\aind_{SO(2)}\rho_2=2\aind_{SO(2)}\lambda_2$。注意 $\aind_{SO(2)}\lambda_2=\chi(S^2)=2$，于是 $\aind_{SO(2)}\lambda_2=1$。

\transnote{原文此处写「$\aind_{SO(2)}\lambda_2=\chi(S^2)=2$」，与上一句 $\aind_{SO(2)}\rho_2=2\aind_{SO(2)}\lambda_2$ 自相矛盾：$\lambda_2$ 是 Thom 类，$\chi(S^2)$ 应由 de Rham 象征 $\rho_2$ 给出。按上下文应为 $\aind_{SO(2)}\rho_2=\chi(S^2)=2$。译文照录原文，读者按此理解。}

$n=1$ 的情形。回忆我们讨论伪微分算子时提到的圆 $S^1$ 上的椭圆算子 $A$。经过适当同伦得到 $-h(\lambda_1)=\sigma(A)$，其中 $h$ 是自然延拓同态 $K(T\R^1)\to K(TS^1)$。由 $\Ind A=-1$ 可知 $\aind_{O(1)}\lambda_1=1$。这个构造在附录里作为粘合函数的例子详细展开。

把 $\R^n$ 分解成若干 $\R^1$ 与 $\R^2$ 的因子。由乘积公式可知，对 $O(n)$ 的每个子群 $G=G_1\times\cdots\times G_n$（各 $G_i$ 取 $O(1)$ 或 $SO(2)$）都有 $\aind_G\lambda_n=1$。$O(n)$ 的表示由它到这些子群的限制完全决定，故 $\aind_{O(n)}\lambda_n=1$。
\end{proof}

设 $\varphi:K(TX)\to K(TY)$ 是复向量丛 $TY\to TX$ 的 Thom 同构。把上面两个命题合起来得到
\[
\begin{aligned}
\aind^Y\circ\varphi(a)
&=\aind^Y(\pi^*a\cdot\lambda_{TY})\\
&=(\aind^X a)\cdot\big(\aind_{O(n)}^{\R^n}\lambda_n\big)=\aind^X a.
\end{aligned}
\]
至此已证解析指标与 Thom 同构交换，Atiyah–Singer 指标定理的证明完毕。

% ===== 原书 p.32 =====
\section{等变 K 理论与齐性空间}
\subsection{G-向量丛与 $K_G(X)$}

下面考虑 K 理论的推广，把紧 Lie 群 $G$ 的作用计入。回忆 Lie 群是带有光滑流形结构的群，其乘法与取逆映射都光滑。特别地，每个有限群都是零维 Lie 群。所谓流形 $X$ 上的左 $G$-作用，是指光滑映射 $G\times X\to X$，记为 $(g,x)\mapsto g\cdot x$，满足 $(gg')\cdot x=g\cdot(g'\cdot x)$ 与 $e\cdot x=x$。类似地，右作用是一个映射 $(g,x)\mapsto x\cdot g$，满足 $x\cdot(gg')=(x\cdot g)\cdot g'$。两种情形下分别称 $G$ 从左侧或右侧作用在 $X$ 上。所谓 $G$-空间，是指流形 $X$ 连同指定的 $G$-作用。

\begin{definition}
称 $X$ 上的向量丛 $E$ 是\emph{G-向量丛}（G-vector bundle），若 $E$ 与 $X$ 都是 $G$-空间，且下面两个条件成立：
\begin{enumerate}[label=\arabic*.,leftmargin=2.2em]
\item 丛投影 $\pi:E\to X$ 与 $E$、$X$ 上的 $G$-作用交换。换言之，对所有 $v\in E$、$g\in G$ 有 $g\cdot\pi(v)=\pi(g\cdot v)$；
\item 对每个 $g\in G$ 与每个 $x\in X$，映射 $E_x\to E_{g\cdot x}$，$v\mapsto g\cdot v$，都是线性的。
\end{enumerate}
\end{definition}

\begin{definition}
给定 $G$-空间 $X$，把 Grothendieck 构造施于 $\Vect_G(X)$，即 $X$ 上复 $G$-向量丛的 $G$-同构类构成的半群，所得群记作 $K_G(X)$。对 $K_G(X)$ 的研究称为\emph{等变 K 理论}。
\end{definition}

§1 中关于普通 K 理论的结论都可直接搬到等变情形，只需把记号作相应改动。最明显的例外有两条：分裂原理只在 $G$ 交换时成立；等变 Thom 同构定理的证明要困难得多。等变 K 理论的完整论述见 [17]。

\begin{definition}
给定连续群同态 $\varphi:H\to G$，任一 $G$-空间 $X$ 都可看作 $H$-空间，作用取 $h\cdot x=\varphi(h)\cdot x$。把这一构造施于 $X$ 上的 $G$-向量丛，即得映射 $\varphi^*:K_G(X)\to K_H(X)$。
\end{definition}

注意若 $G$ 是平凡群，则 $K_G(X)$ 化归普通 K 理论中的环 $K(X)$。又注意若 $X$ 是一点，则 $X$ 上的复 $G$-向量丛就是复 $G$-模。此时 $K_G(\pt)$ 成为 $G$ 的表示环 $R(G)$，它由把 Grothendieck 构造施于 $G$ 的有限维复表示空间构成的半群得到。一般地，考虑由 $X$ 到一点的映射诱导的同态 $R(G)\to K_G(X)$，可知 $K_G(X)$ 是 $R(G)$-模。一个极端情形是

\begin{proposition}
若 $G$ 平凡地作用在 $X$ 上（即 $g\cdot x=x$ 对一切 $g$），则 $K_G(X)\cong K(X)\otimes R(G)$。
\end{proposition}

证明见 [17, §2.2]；有限群情形见 [6, §1.6]。另一个极端给出

\begin{proposition}
若 $G$ 自由地作用在 $X$ 上（即 $g\cdot x=x$ 当且仅当 $g=e$），则 $K_G(X)\cong K(X/G)$，其中 $X/G$ 是把 $x$ 与所有 $g\cdot x$ 等同所得的轨道空间（orbit space）。
\end{proposition}

\begin{proof}
设 $E$ 是 $X$ 上的 $G$-向量丛。$G$-作用此时平凡，故 $E/G$ 是 $X/G$ 上的普通向量丛。\footnote{还须验证 $E/G$ 局部平凡。}于是得到同态 $K_G(X)\to K(X/G)$。

% ===== 原书 p.33 =====
反过来，考虑由标准投影 $p:X\to X/G$ 诱导的映射。给定 $X/G$ 上的普通向量丛 $F$，拉回丛 $p^*F$ 的全空间为
\[
\{(x,v)\in X\times F\mid p(x)=\pi(v)\}.
\]
赋予 $G$-作用 $g\cdot(x,v)=(g\cdot x,v)$，可把 $p^*F$ 做成 $X$ 上的 $G$-向量丛，从而得到反向同态 $K(X/G)\to K_G(X)$。由典范向量丛同构 $E\to p^*(E/G)$ 与 $(p^*F)/G\to F$，结论立即成立。
\end{proof}

\transnote{本证明首句原文作「$G$-作用此时平凡」，与命题的自由作用前提相悖，应读作「自由」；正是自由性让 $E/G$ 成为 $X/G$ 上的普通向量丛。另外，同构 $K_G(X)\cong K(X)\otimes R(G)$ 只在作用平凡时成立，此处 $\otimes$ 是普通张量积，原文未加帽。}

\subsection{齐性空间}

从等变 K 理论的视角看，特别有趣的一类流形是齐性空间。设 $H$ 是 $G$ 的闭子群，考虑左陪集商空间 $G/H$，赋以商拓扑。所得空间是光滑流形（见 [1, p. 21]）。若 $G$ 传递地作用在流形 $X$ 上（transitive action，即对任意 $x,y\in X$，存在 $g\in G$ 使 $g\cdot x=y$），则称 $X$ 是\emph{齐性空间}。进一步，取 $H$ 为稳定子群（即固定某点的子群），则有同构 $X\cong G/H$。

齐性空间的一些简单例子如下：
\begin{itemize}[leftmargin=1.6em]
\item 任一 Lie 群：$G\cong G/\{e\}$；
\item 球面：$S^n\cong\SO(n+1)/\SO(n)$；
\item 奇维球面（odd sphere）：$S^{2n+1}\cong\U(n+1)/\U(n)$，或 $S^{2n+1}\cong\SU(n+1)/\SU(n)$（$n>0$）；
\item 实 Grassmann 流形：$G(n,k)\cong O(n+k)/O(n)\times O(k)$；
\item 实射影空间：$\R P^n\cong O(n+1)/O(n)\times O(1)$；
\item 复 Grassmann 流形：$G_{\C}(n,k)\cong\U(n+k)/\U(n)\times\U(k)$；
\item 复射影空间：$\C P^n\cong\U(n+1)/\U(n)\times\U(1)$；
\item 透镜空间（lens space）：$L(n,k)=S^{2n+1}/\Z_k\cong\U(n+1)/\U(n)\times\Z_k$。
\end{itemize}

由上述命题，注意 $H$ 自由地作用在 $G$ 上（标准作用 $h\cdot g=hg$），得到 $K(G/H)\cong K_H(G)$。不过我们的兴趣更在 $G/H$ 上的 $G$-向量丛。给定有限维 $H$-模 $M$，可构造乘积 $G\times_H M$，即 $G\times M$ 模去 $H$ 的作用：
\[
G\times_H M=G\times M\Big/(g,x)\sim h\cdot(g,x)=(gh^{-1},hx).
\]
设 $\pi:G\times_H M\to G/H$ 是投影映射 $(g,x)\mapsto gH$。它与 $G/H$ 及 $G\times_H M$ 上的标准 $G$-作用 $g'\cdot gH=g'gH$、$g'\cdot(g,x)=(g'g,x)$ 显然交换，因此 $G\times_H M$ 是 $G/H$ 上以 $M$ 为纤维的向量丛。事实上下面证明，$G/H$ 上任一 $G$-向量丛都可由这一构造得到。

\begin{proposition}
若 $H$ 是 Lie 群 $G$ 的闭子群，则 $K_G(G/H)\cong R(H)$。
\end{proposition}

\begin{proof}
设有 $G/H$ 上的 $G$-向量丛 $E$。$E$ 上的 $G$-作用限制为 $E_H$ 上的 $H$-作用，故 $E$ 在陪集 $H$ 处的纤维是有限维 $H$-模 $M$。$E$ 上的 $G$-作用 $G\times E\to E$ 限制为 $G$-同态 $G\times M\to E$，它在 $G\times M$ 的 $H$-作用下不变，故诱导 $G$-映射 $\alpha:G\times_H M\to E$。要证 $\alpha$ 是 $G$-同构。由 $G$-向量丛的定义，映射

% ===== 原书 p.34 =====
$E_H\to E_{gH}$，$x\mapsto g\cdot x$，对每个 $g\in G$ 都是同构（其逆为 $x\mapsto g^{-1}\cdot x$）。把 $\alpha$ 限制到任一陪集 $gH$ 处的纤维，得 $\alpha|_{gH}:gH\times_H M\to E_{gH}$，它恰是上述同构。于是 $\alpha$ 是整体的向量丛同构。
\end{proof}

\transnote{这一命题是 §4 的核心机制：$G$ 传递地作用在 $G/H$ 上，丛的整体行为由它在单点 $H$ 处的纤维（一个 $H$-模）决定，所以等变 K 理论把 $G/H$ 上的几何问题化归为 $H$ 的表示论。后文的 Bott 定理正是沿这条思路把指标化成表示环里的元素。}

\begin{example}
考虑球面 $S^2\cong\SO(3)/\SO(2)$，得
\[
K_{\SO(3)}(S^2)\cong R(\SO(2))\cong\Z[z,z^{-1}].
\]
特别地，这个环比 $K(S^2)\cong\Z^2$ 大得多。
\end{example}

用 $G$-作用在齐性空间 $G/H$ 上「平移」，可以看出等变向量丛的整体行为由它在单点的行为决定。用完全类似的论证，可证上述结果的微小推广。

\begin{proposition}
若 $X$ 是局部紧\footnote{这里用紧支撑 $K_G$ 理论，定义方式与普通情形类似。}的 Hausdorff $H$-空间，则 $K_G(G\times_H X)\cong K_H(X)$。
\end{proposition}

\begin{example}
齐性空间 $G/H$ 的切丛是实向量丛，在陪集 $H$ 处的纤维为 $\mathfrak{g}/\mathfrak{h}$，其中 $\mathfrak{g}$、$\mathfrak{h}$ 分别是 $G$、$H$ 的 Lie 代数。\footnote{Lie 群 $G$ 的 Lie 代数 $\mathfrak{g}\cong TG_e$ 是 $G$ 在单位元处的切空间，括号由左不变向量场的括号给出。}注意 $T(G/H)$ 是 $G$-向量丛，且 $\mathfrak{g}/\mathfrak{h}$ 带有自然的 $H$-作用（由映射 $A_h:gH\mapsto hgH$ 诱导的作用 $\mathrm{Ad}:H\to\Aut(\mathfrak{g}/\mathfrak{h})$），可知 $T(G/H)\cong G\times_H(\mathfrak{g}/\mathfrak{h})$。于是得到同构
\[
K_G\big(T(G/H)\big)\cong K_G\big(G\times_H(\mathfrak{g}/\mathfrak{h})\big)\cong K_H(\mathfrak{g}/\mathfrak{h}).
\]
\end{example}

\subsection{G-指标定理}

若 $E$ 是紧 $G$-空间 $X$ 上的 $G$-向量丛，则可把光滑截面空间 $\Gamma E$ 做成（无限维）$G$-模：

\begin{definition}
若 $s:X\to E$ 是光滑截面（不必与 $X$、$E$ 上的 $G$-作用交换），则给 $\Gamma E$ 赋予 $G$-作用 $(g\cdot s)(x)=g\cdot s(g^{-1}\cdot x)$。
\end{definition}

给定 $X$ 上的复 $G$-向量丛 $E$、$F$，自然要考察与 $\Gamma E$、$\Gamma F$ 上 $G$-作用交换的椭圆算子 $P:\Gamma E\to\Gamma F$。此时称 $P$ 是\emph{G-不变的}。注意 $TX$ 是 $G$-空间\footnote{设 $L_g:X\to X$ 是映射 $x\mapsto g\cdot x$，$dL_g:TX\to TX$ 是它的微分。取 $g\cdot(x,\xi)=dL_g(x,\xi)$，即得 $TX$ 上的标准 $G$-作用。}，$\pi^*E$、$\pi^*F$ 都是 $TX$ 上的 $G$-向量丛，故象征 $\sigma(P):\pi^*E\to\pi^*F$ 是 $G$-映射，从而给出 $K_G(TX)$ 中的一个元素。与前面一样，等变情形下的拓扑指标定义为合成
\[
\tind_G^X:K_G(TX)\xrightarrow{\ i_!\ }K_G(TM)\xrightarrow{\ j_!^{-1}\ }K_G(\pt)\cong R(G),
\]
其中 $i:X\to M$ 是把 $X$ 嵌入有限维实 $G$-模 $M$ 的嵌入，$j:\pt\to M$ 是原点的包含。

由于 $P$ 是椭圆的，$\Ker P$ 与 $\Coker P$ 都是有限维的。又因 $P$ 是 $G$-不变的，$\Ker P$、$\Coker P$ 在 $\Gamma E$、$\Gamma F$ 的 $G$-作用下不变。于是 $\Gamma E$、$\Gamma F$ 上的 $G$-作用诱导 $\Ker P$、$\Coker P$ 上的 $G$-作用，$\Ker P$、$\Coker P$ 都是有限维 $G$-模。据此可定义

% ===== 原书 p.35 =====
\begin{definition}
$G$-不变椭圆算子 $P$ 的 $G$-指标\footnote{$G$-指标与普通指标由等式 $\Ind P=\dim\Ind_G P$ 相联系。}为
\[
\Ind_G=[\Ker P]-[\Coker P]\in R(G).
\]
\end{definition}

利用 $G$-指标可把解析指标定义为同态
\[
\aind_G^X:K_G(TX)\to K_G(\pt)\cong R(G),
\]
由 $\aind_G^X=\Ind_G P$ 给出，其中 $P$ 是 $G$-不变椭圆算子且 $\sigma(P)=x$。§3 中给出的 Atiyah–Singer 指标定理证明可自然地搬到等变情形（作相应的必要改动：把 $\Z$ 换成 $R(G)$，各处补上 $G$）。于是得到 $\aind_G=\tind_G$，G-指标定理随即成立。

\begin{theorem}[G-指标定理]
若 $P:\Gamma E\to\Gamma F$ 是紧流形 $X$ 上的 $G$-不变椭圆（伪微分）算子，则
\[
\Ind_G P=\tind_G^X\sigma(P).
\]
\end{theorem}

\subsection{诱导表示}

本节回顾 Lie 群表示论的一些基本结论，用表示环的记号叙述。更完整的论述见 [1]。以下取 $G$ 为紧 Lie 群。称 $G$-模是\emph{不可约的}，若它没有 $G$-不变真子空间；称两个 $G$-模 $M$、$N$ \emph{等价}，若存在 $G$-同构 $M\to N$。取 $\{w_\alpha\}$ 为 $G$ 的一族完全的不等价不可约表示，则表示环 $R(G)$ 是由 $\{w_\alpha\}$ 生成的自由阿贝尔群。定义两个 $G$-模的\emph{缠绕数}（intertwining number）为
\[
\tau(M,N)=\dim_\C\Hom_G(M,N),
\]
可把它扩充为 $R(G)$ 上的内积 $\ip{\cdot}{\cdot}_G$。由 Schur 引理，若 $M$、$N$ 都不可约，则
\[
\ip{M}{N}_G=
\begin{cases}
1, & M\text{ 与 }N\text{ 等价},\\
0, & \text{其他},
\end{cases}
\]
于是 $\{w_\alpha\}$ 关于此内积正交归一（orthonormal）。

其次，形式地定义表示环的无限扩张（infinite extension）$\widehat R(G)$，它是所有可能无限的线性组合 $\sum_\alpha a_\alpha w_\alpha$（$a_\alpha\in\Z$）构成的集合。显然有自然包含 $R(G)\subset\widehat R(G)$。$R(G)$ 上的内积随即扩充为配对（pairing）$R(G)\otimes\widehat R(G)\to\Z$：
\[
\Big\langle x,\sum_\alpha a_\alpha w_\alpha\Big\rangle_G=\sum_\alpha a_\alpha\ip{x}{w_\alpha}_G,
\]
其中 $x\in R(G)$，$\sum_\alpha a_\alpha w_\alpha\in\widehat R(G)$。注意此和其实是有限和，因为除有限多个 $\alpha$ 外 $\ip{x}{w_\alpha}$ 为零。

\begin{definition}
给定连续群同态 $\varphi:H\to G$，定义形式诱导表示（formal induced representation）$\varphi_*:R(H)\to\widehat R(G)$ 为映射
\[
\varphi_*:x\mapsto\sum_\alpha\ip{\varphi^*w_\alpha}{x}_H w_\alpha,
\]
其中 $\varphi^*:R(G)\to R(H)$ 是等变 K 理论诱导的映射。
\end{definition}

\transnote{求和指标 $\alpha$ 取遍 $G$ 的不可约表示，而配对 $\ip{\varphi^*w_\alpha}{x}_H$ 在 $H$ 上进行：$x\in R(H)$，$\varphi^*w_\alpha$ 是 $w_\alpha$ 限制成 $H$-模所得。原文紧接着的伴随关系把 $x$ 印作 $x\in R(G)$，应为 $R(H)$。}

% ===== 原书 p.36 =====
把上述定义用扩充内积展开，可知诱导表示满足伴随关系
\[
\ip{y}{\varphi_*x}_G=\ip{\varphi^*y}{x}_H
\]
对所有 $x\in R(G)$、$y\in R(G)$ 成立。若 $H$ 是 $G$ 的闭子群，$i:H\to G$ 是包含映射，则可为诱导表示构造一个几何上更直观的替代定义。特别地，给定有限维 $H$-模 $M$，回忆 $G/H$ 上向量丛 $G\times_H M$ 的光滑截面空间 $\Gamma(G\times_H M)$ 具有无限维 $G$-模结构。把这一映射延拓到表示环，得到

\begin{definition}
给定闭子群 $H$ 的包含映射 $i:H\to G$，定义\emph{诱导表示} $i_*:R(H)\to\widehat R(G)$ 为映射
\[
M\mapsto\Gamma(G\times_H M)
\]
到 $R(H)$ 的延拓，其中 $M$ 是有限维 $H$-模。
\end{definition}

为说明这两个定义等价，需要证明下面的 Frobenius 定理：

\begin{theorem}[Frobenius 互反]
若 $W$ 是 $G$-模，$M$ 是 $H$-模，则
\[
\Hom_G(W,i_*M)\cong\Hom_H(i^*W,M),
\]
其中 $i_*M=\Gamma(G\times_H M)$，$i^*W$ 是 $W$ 限制为 $H$-模所得。
\end{theorem}

\begin{proof}
设有 $G$-不变同态 $F:W\to\Gamma(G\times_H M)$。把它与映射 $\Gamma(G\times_H M)\to M$，$s\mapsto s(H)$，复合，得到相应的 $H$-同态 $f:W\to M$。注意对 $h\in H$、$w\in W$ 有
\[
\begin{aligned}
f(h\cdot w)&=F(h\cdot w)(H)=(h\cdot F(W))(H)\\
&=h\cdot\big(F(w)(h^{-1}H)\big)=h\cdot\big(F(w)(H)\big)=h\cdot f(w).
\end{aligned}
\]
反过来，给定 $H$-不变同态 $f:W\to M$，考虑相应的映射 $F:W\to\Gamma(G\times_H M)$，它对 $gH\in G/H$、$w\in W$ 定义为
\[
F(w)(gH)=\big(g,f(g^{-1}\cdot w)\big)\in G\times_H M.
\]
由于 $f:W\to M$ 是 $H$-不变的，此映射良定义。此外注意
\[
\begin{aligned}
F(g'\cdot w)(gH)&=\big(g,f(g^{-1}g'\cdot w)\big)
=g'\cdot\Big(g'^{-1}g,f\big((g'^{-1}g)^{-1}\cdot w\big)\Big)\\
&=g'\cdot F(w)(g'^{-1}\cdot gH)=\big(g'\cdot F(w)\big)(gH)
\end{aligned}
\]
对所有 $g\in G$ 成立。故 $F:W\to\Gamma(G\times_H M)$ 是 $G$-同态。注意两个映射 $f\mapsto F$ 与 $F\mapsto f$ 互为逆，结论立即成立。
\end{proof}

\begin{corollary}
若 $i:H\to G$ 是闭子群 $H$ 在 $G$ 中的包含映射，则对任一 $H$-模 $M$ 有
\[
\big[\Gamma(G\times_H M)\big]=\sum_\alpha\ip{i^*w_\alpha}{[M]}_H w_\alpha\in\widehat R(G),
\]
于是诱导表示的两个定义一致。
\end{corollary}

\begin{proof}
由 Frobenius 互反定理，
\[
\begin{aligned}
\big[\Gamma(G\times_H M)\big]
&=\sum_\alpha\dim\Hom_G\big(w_\alpha,\Gamma(G\times_H M)\big)\cdot w_\alpha\\
&=\sum_\alpha\dim\Hom_H(i^*w_\alpha,M)\cdot w_\alpha
=\sum_\alpha\ip{i^*w_\alpha}{[M]}_H w_\alpha,
\end{aligned}
\]
这正是所要。
\end{proof}

% ===== 原书 p.37 =====
\subsection{齐性微分算子}

作为 G-指标定理的特例，下面考察齐性空间 $G/H$ 上的 $G$-不变伪微分算子\footnote{本节所有结论对伪微分算子同样成立，但为记号简便，以下省略「伪」字。} $P:\Gamma E\to\Gamma F$ 的 $G$-指标。此时称 $P$ 是\emph{齐性微分算子}。借助 Peter-Weyl 定理，此情形的分析归结为诱导映射的表示论。又因齐性算子这类对象附加的结构过多，下面很快会看到，$P$ 的指标只依赖向量丛 $E$、$F$，并不直接依赖算子本身。仿照 [12]，定义齐性象征如下：

\begin{definition}
对 $G/H$ 上的齐性算子 $P:\Gamma E\to\Gamma F$，回忆可写 $E\cong G\times_H M$、$F\cong G\times_H N$，其中 $M$、$N$ 是有限维 $H$-模。于是定义 $P$ 的\emph{齐性象征}（homogeneous symbol）为
\[
\hat\sigma(P)=[M]-[N]\in R(H).
\]
\end{definition}

现在把齐性象征与普通象征 $\sigma(P)\in K_G(T(G/H))$ 作比较。后者由 $T(G/H)$ 上的齐性复形
\[
0\longrightarrow\pi^*E\xrightarrow{\ \sigma\ }\pi^*F\longrightarrow 0
\]
表示，其中 $\sigma:\pi^*E\to\pi^*F$ 是 $G$-同态。设 $s:\pt\to T(G/H)$ 表示把恒等陪集 $H\in G/H$ 经零截面映入 $T(G/H)$ 的包含映射。诱导映射 $s^*:K_G(T(G/H))\to K_G(\pt)=R(G)$ 把 $\sigma(P)$ 映为一点上的复形\footnote{注意映射 $\sigma(H,0):M\to N$ 是同态，但一般不是 $H$-同态。}
\[
0\longrightarrow M\xrightarrow{\ \sigma(H,0)\ }N\longrightarrow 0.
\]
于是
\[
s^*\sigma(P)=[M]-[N]=\hat\sigma(P)\in R(H).
\]
现在可计算齐性微分算子的 $G$-指标。

\begin{theorem}[Bott]
若 $P$ 是 $G/H$ 上的齐性微分算子，则
\[
\Ind_G P=i_*\hat\sigma(P)=i_*\circ s^*\sigma(P)\in\widehat R(G),
\]
其中 $i_*:R(H)\to\widehat R(G)$ 是形式诱导同态。特别地，$i_*\hat\sigma(P)$ 其实是 $R(G)\subset\widehat R(G)$ 中的有限元素。
\end{theorem}

\transnote{等号两端所在的环不同，读时须小心：$s^*\sigma(P)$ 落在 $R(G)$，而 $[M]-[N]$ 是 $H$-模的形式差，落在 $R(H)$。原因在于 $s$ 把点送到 $(H,0)$，该点的迷向群是 $H$ 而非 $G$；$s$ 本身不是 $G$-映射，故 $s^*\sigma(P)$ 只带 $H$-对称性。}

由诱导表示的替代定义，
\[
i_*\big(\hat\sigma(P)\big)=i_*\big([M]-[N]\big)=\big[\Gamma(G\times_H M)\big]-\big[\Gamma(G\times_H N)\big]=[\Gamma E]-[\Gamma F],
\]
于是 Bott 定理可改写成更简洁的形式。

\begin{theorem}
若 $P:\Gamma E\to\Gamma F$ 是齐性微分算子，则
\[
\Ind_G P=[\Gamma E]-[\Gamma F]\in\widehat R(G),
\]
其中 $[\Gamma E]-[\Gamma F]$ 其实是 $R(G)\subset\widehat R(G)$ 中的有限元素。
\end{theorem}

% ===== 原书 p.38 =====
为证明此定理，须把无限维 $G$-模 $\Gamma E$ 分解成有限维 $G$-不变子空间。限制在这些子空间上，定理容易验证，再把结论提升到光滑截面的整体空间。这些子空间的自然候选如下。

\begin{definition}
对每个 $w_\alpha$，定义 $\Gamma_\alpha E$ 为同态
\[
i_\alpha:w_\alpha\otimes\Hom_G(w_\alpha,\Gamma E)\to\Gamma E
\]
的像，其中 $i_\alpha:(w\otimes\varphi)\mapsto\varphi(w)$。于是 $\Gamma_\alpha(E)$ 是 $\Gamma E$ 的有限维 $G$-不变子空间，其在表示环中的像为 $[\Gamma_\alpha E]=\ip{w_\alpha}{[\Gamma E]}_G w_\alpha\in R(G)$。
\end{definition}

注意 $\Gamma E$ 在表示环的无限扩张 $\widehat R(G)$ 中的类由 $[\Gamma E]=\sum_\alpha\ip{w_\alpha}{[\Gamma E]}_G w_\alpha$ 给出。此式提示可（至少形式地）得到无限分解 $\Gamma E=\oplus_\alpha\Gamma_\alpha E$。下述著名结果把这一分解落实。

\begin{lemma}[Peter-Weyl]
设 $L^2(E)$ 是把 $\Gamma E$ 关于某个 $G$-不变内积完备化所得的 Hilbert 空间，则诸子空间 $\Gamma_\alpha E$ 构成 $L^2(E)$ 的一族完备正交子空间。
\end{lemma}

现在已有足够准备，可较快证明 Bott 定理。

\begin{proof}[Bott 定理的证明]
设 $P:\Gamma E\to\Gamma F$ 是椭圆齐性微分算子。因 $P$ 是 $G$-映射，它在每个有限维 $G$-不变子空间 $\Gamma_\alpha$ 上诱导线性映射 $P_\alpha:\Gamma_\alpha E\to\Gamma_\alpha F$。因 $P$ 椭圆，$\Ker P$ 与 $\Coker P$ 是有限维 $G$-模，Peter-Weyl 定理给出
\[
\Ker P=\bigoplus_\alpha\Ker P_\alpha,\qquad
\Coker P=\bigoplus_\alpha\Coker P_\alpha,
\]
其中除有限多个 $\alpha$ 外 $\Ker P_\alpha$ 与 $\Coker P_\alpha$ 为零。于是
\[
\Ind_G P=[\Ker P]-[\Coker P]=\sum_\alpha[\Ker P_\alpha]-[\Coker P_\alpha].
\]
又注意有有限维 $G$-模的正合列
\[
0\longrightarrow\Ker P_\alpha\longrightarrow\Gamma_\alpha E
\xrightarrow{\ P_\alpha\ }\Gamma_\alpha F\longrightarrow\Coker P_\alpha\longrightarrow 0,
\]
可对 $R(G)$ 使用由复形给出的 $K_G$ 理论构造。此复形的支撑为空，其欧拉示性数（即各 $G$-模的交错和）在 $R(G)$ 中必为零。于是
\[
\begin{aligned}
[\Ker P_\alpha]-[\Coker P_\alpha]
&=[\Gamma_\alpha E]-[\Gamma_\alpha F]\\
&=\ip{w_\alpha}{[\Gamma E]}_G w_\alpha=\ip{w_\alpha}{[\Gamma F]}_G w_\alpha.
\end{aligned}
\]
对 $\alpha$ 求和，得到
\[
\Ind_G P=\sum_\alpha\ip{w_\alpha}{[\Gamma E]}_G w_\alpha-\sum_\alpha\ip{w_\alpha}{[\Gamma F]}_G w_\alpha=[\Gamma E]-[\Gamma F],
\]
这正是所要的结论。
\end{proof}

\transnote{「支撑为空」是说该复形的象征在无穷远处可逆，故复形在 $K_G$ 中代表零类；等价地，它在 $R(G)$ 中的欧拉示性数 $\sum_i(-1)^i[\Gamma_\alpha^i]$ 为零。正是这一点把 $[\Ker P_\alpha]-[\Coker P_\alpha]$ 换成 $[\Gamma_\alpha E]-[\Gamma_\alpha F]$。}


% ---- 附录：改用字母编号（原书 A、B） ----
\setcounter{section}{0}
\renewcommand{\thesection}{\Alph{section}}
\renewcommand{\seclabel}{附录 \thesection}
% ===== 原书 p.39 =====
% 注：p.39 上半页属 §4 末例（圆的例子），按节界归 c06；本块自 \section 起。
\section{Bott 周期性的一个证明}
\subsection{预备}

正文始终用上同调、K 理论与表示论这套工具，从分析一侧求解指标问题。本附录把方向反过来：把指标理论用到拓扑上。具体做法是用一族椭圆算子的指标，给 Bott 周期性定理一个简单证明。§1 已说明，Bott 周期性定理是 K 理论的基本基石，使它具有周期上同调理论的结构。希望这个证明能让读者看清代数拓扑与椭圆算子之间的根本联系（若正文其余部分还没能让读者看清）。这一联系的更多材料见 [3]。我们的证明沿 [5] 的思路，是 [4] 中更一般论证的一个稍简情形。本附录后段将简述：如何把这一般论证推广，去证明 Thom 同构定理的等变版本。

开始前先回顾 §1 中 Bott 周期性定理的陈述。设 $b$ 是 $\tilde K(S^2)=K(\C)$ 的生成元。注意外张量积（exterior tensor product）$\BT:K(X)\otimes K(Y)\to K(X\times Y)$ 限制为乘积 $\BT:\tilde K(X)\otimes\tilde K(Y)\to\tilde K(X\wedge Y)$，\footnote{利用分解 $K(X)\iso\tilde K(X)\oplus K(\pt)$，外张量积分裂为
\[
\otimes:\bigl(\tilde K(X)\otimes\tilde K(Y)\bigr)\oplus\bigl(\tilde K(X)\oplus\tilde K(Y)\bigr)\oplus K(\pt)\longrightarrow\tilde K(X\wedge Y)\oplus\tilde K(X\vee Y)\oplus K(\pt).
\]
由 $\tilde K(X)\oplus\tilde K(Y)\to\tilde K(X\vee Y)$ 可知，外张量积限制为乘积映射
\[
\BT:\tilde K(X)\otimes\tilde K(Y)\to\tilde K(S\wedge Y),
\]
其中回忆 $X\vee Y=(X\times\pt)\cup(\pt\times Y)$ 是单点并，$X\wedge Y=(X\times Y)/(X\vee Y)$ 是 smash 积。} 据此定义同态
\[
\beta:K(X)\to K^{-2}(X)=\tilde K(S^2\wedge X^+).
\]

\transnote{本附录原文有若干处符号笔误，译文一律照原样排印。已知的有：上面脚注中的 $\tilde K(S\wedge Y)$ 应为 $\tilde K(X\wedge Y)$；下一段开头的 $\beta(X)$ 应为 $\beta(x)$；p.40 的 $\sum_{0\le m<n}b_m e^{imx}$ 应为 $e^{im\theta}$；p.42 的 $f(x)=z^{-2}$ 应为 $f(z)$；p.44 的 $\alpha(X)\cdot y$ 应为 $\alpha(x)\cdot y$，同页 $u\cdot\lambda_1$ 应为 $u\cdot b$，$K(X\times\R^n)$ 应为 $K(X\times\R^2)$；同页 $\alpha\beta(X)$ 应为 $\alpha\beta(x)$。读到这些位置按上下文理解即可。}

% ===== 原书 p.40 =====
映射由 $\beta(X)=b\BT x$ 给出。Bott 周期性定理断言这个映射是同构。用紧支撑 K 理论注意到
\[
K(\R^n\times X)=\tilde K\bigl((\R^n\times X)^+\bigr)=\tilde K(S^n\wedge X^+)=K^{-n}(X),
\]
于是可以把 $\beta$ 看成同态
\[
\beta:K(X)\xrightarrow{\,b\BT\cdot\,}K(\R^2\times X),
\]
其中 $b$ 视为 $K(\R^2)$ 的元素。为证明 Bott 周期性定理，我们将用一族椭圆算子的指标，构造逆同态 $\alpha:K(\R^2\times X)\to K(X)$。

\subsection{Wiener--Hopf 算子}

先看圆周 $S^1=\R/2\pi\Z$ 上一个具体的伪微分算子。§3 引入过 $S^1$ 上复值函数的零阶伪微分算子 $P$，它在 Fourier 级数下由
\[
P(e^{ik\theta})=\begin{cases} e^{ik\theta} & k\ge 0,\\ 0 & k<0 \end{cases}
\]
给出。这个算子可以看成从 $S^1$ 上复值函数空间到「只有正 Fourier 系数」的子空间的投影。设 $f:S^1\to\C^*$ 连续，其中 $\C^*$ 是非零复数。可以构造相应的伪微分算子 $P_f=PfP+(1-P)$。写 $f(\theta)=\sum_n a_n e^{in\theta}$，按 Fourier 级数展开得到
\[
P_f(e^{ik\theta})=\begin{cases} \sum_{n>-k}a_n e^{i(n+k)\theta} & k\ge 0,\\ e^{ik\theta} & k<0, \end{cases}
\]
注意 $P_f$ 是椭圆的，因为它的零阶象征由
\[
\sigma_{P_f}(\theta,\xi)=\begin{cases} f(\theta) & \xi>0,\\ 1 & \xi<0 \end{cases}
\]
给出。

\transnote{$P$ 的像含常数项：上图对 $k=0$ 也取 $e^{ik\theta}$，故原文「只有正 Fourier 系数」应读作「非负」。这也暴露下一式的下界问题：按 $P_f=PfP+(1-P)$ 逐项计算，求和应为 $\sum_{n\ge-k}$；原文写 $\sum_{n>-k}$，在 $k=0$ 时会漏掉 $a_0$，而原文随后断言 $\Coker P_n$ 为 $n$ 维，只有用 $\sum_{n\ge-k}$ 才对得上。此处按上下文补边界项理解。}

特别地，令 $P_n=P_f$，其中 $f(\theta)=e^{in\theta}$。当 $n\ge 0$ 时有
\[
P_n(e^{ik\theta})=\begin{cases} e^{i(n+k)\theta} & k\ge 0,\\ e^{ik\theta} & k<0, \end{cases}
\]
于是 $\Ker P_n$ 平凡，而 $\Coker P_n$ 是 $n$ 维空间，由形如 $g(\theta)=\sum_{0\le m<n}b_m e^{imx}$ 的函数组成。故 $\Ind P_n=-n$。同理，当 $n<0$ 时有
\[
P_n(e^{ik\theta})=\begin{cases} e^{i(n+k)\theta} & k\ge -n,\\ 0 & 0\le k<-n,\\ e^{ik\theta} & k<0. \end{cases}
\]
此时 $\Coker P_n$ 平凡，而 $\Ker P_n\iso\Coker P_{-n}$，于是同样得到 $\Ind P_n=-n$。

现在回顾：任何连续映射 $f:S^1\to\C^*$ 都同伦于形如 $f(\theta)=e^{in\theta}$ 的映射，映射的度定义为 $\deg f=n$。% ===== 原书 p.41 =====
指标是局部常值函数，于是 $\Ind P_f=-\deg f$。

现在推广这一构造。第一步，不妨把讨论扩充到 $S^1$ 上复向量值函数。设 $f:S^1\to\GL(N,\C)$ 连续，可以构造相应的椭圆算子 $P_f$。注意 $\Ind P_f$ 只依赖 $f$ 的同伦类；这里同伦类就是 $\GL(N,\C)$ 基本群中的类 $[f]\in\pi_1(\GL(N,\C))$。再进一步：设 $V$ 是紧空间 $X$ 上的复向量丛，并对每个 $x\in X$ 给出一族连续变化的连续映射 $f_x:S^1\to\Aut V_x$。于是得到 $S^1$ 上一族由 $X$ 参数化的伪微分算子 $P_f$，其中每个 $(P_f)_x=P_{f_x}$ 都椭圆。这样的集合称为椭圆族（elliptic family）。若 $\dim\Ker(P_f)_x$ 与 $\dim\Coker(P_f)_x$ 对所有 $x\in X$ 都是常数，则可把 $\Ker P_f$ 与 $\Coker P_f$ 定义为 $X$ 上相应的向量丛，从而定义
\[
\Ind P_f=[\Ker P_f]-[\Coker P_f]\in K(X).
\]
当 $\dim\Ker(P_f)_x$ 与 $\dim\Coker(P_f)_x$ 不是常数时，Fredholm 算子族的指标的一般构造见 [5, §2]。注意这个指标只依赖映射 $f:(x,\theta)\mapsto\Aut V_x$ 的同伦类 $[f]$。

\subsection{粘合构造}

设 $V$ 是 $X$ 上的向量丛，$f:(x,\theta)\mapsto\Aut V_x$ 是连续映射。如下构造 $S^2\times X$ 上的向量丛 $E$。设 $B^+$ 与 $B^-$ 是球面 $S^2$ 的闭上半球与闭下半球。二者之并 $B^+\cup B^-=S^2$ 是整个球面，交 $B^+\cap B^-=S^1$ 是赤道。取沿映射 $f$ 的并
\[
E=(B^-\times V)\cup_f(B^+\times V),
\]
其中沿交 $(B^-\times V)\cap(B^+\times V)=S^1\times V$ 作等同
\[
(\theta,v)\in B^-\times V_x\ \sim\ \bigl(\theta,f(\theta,x)v\bigr)\in B^+\times V_x.
\]
这个构造称为粘合构造（clutching construction），映射 $f:(x,\theta)\mapsto\Aut V_x$ 称为粘合函数（clutching function）。粘合函数的完整讨论见 [6, §1.4]。

注意 $S^2\times X$ 上每个向量丛 $E$ 都可由这一构造从 $X$ 上某个向量丛 $V$ 与适当的粘合函数 $f:(x,\theta)\mapsto\Aut V_x$ 得到\footnote{设 $i_+,i_-:X\to S^2\times X$ 是 $S^2$ 两个极点处的包含映射。注意 $B^+$ 与 $B^-$ 可缩，故 $E|_{B^+\times X}\iso B^+\times i_+^*E$，且 $E|_{B^-\times X}\iso B^+\times i_-^*E$。取 $V\iso i_-^*E\iso i_+^*E$，再把 $f:(x,\theta)\mapsto\Aut V_x$ 取为 $E|_{B^+\times X}$ 与 $E|_{B^-\times X}$ 在 $S^1\times X$ 上的等同。}；并且 $E$ 由 $f$ 的同伦类 $[f]$ 唯一决定（同构意义下，见 [6, §1.4.6]）。于是得到映射 $\Vect(S^2\times X)\to K(X)$：
\[
E\mapsto (V,[f])\mapsto P_f\mapsto\Ind P_f\in K(X),
\]
它再扩充为群同态 $\alpha:K(S^2\times X)\to K(X)$。注意\footnote{考虑空间序列 $X^+\vee Y^+\to X^+\times Y^+\to X^+\wedge Y^+=(X\times Y)^+$，得到正合列 $0\to K(X\times Y)\to\tilde K(X^+\times Y^+)\to K(X)\oplus K(Y)\to0$。特别地，取 $X=\R^n$、$Y=X$，有正合列 $0\to K(\R^n\times X)\to\tilde K(S^n\times X^+)\to\tilde K(S^n)\oplus K(X)\to0$。注意 $\tilde K(S^n\times X^+)=\tilde K(S^n)\oplus K(S^n\times X)$，可见 $K(S^n\times X)\iso K(\R^n\times X)\oplus K(X)$。} $K(S^2\times X)\iso K(\R^2\times X)\oplus K(X)$，取 $\alpha$ 在 $K(\R^2\times X)$ 上的限制，得到
\[
\alpha:K(\R^2\times X)\to K(X).
\]

\transnote{为什么先落到 $K(S^2\times X)$ 再取出 $K(\R^2\times X)$：粘合构造直接给出的丛住在 $S^2\times X$ 上，而 $K(S^2\times X)\iso K(\R^2\times X)\oplus K(X)$ 的两块按沿 $S^2$ 方向是否平凡分开——$K(X)$ 一块由 $X$ 上的丛沿 $S^2$ 平凡扩张而来，$K(\R^2\times X)$ 一块才携带沿 $S^2$ 的粘合信息。把 $\alpha$ 限制到后者，等于只看在 $S^2$ 方向非平凡的丛。Bott 映射的逆正是在这个直和项上构造的。}

% ===== 原书 p.42 =====
取 $X$ 为一点，则 $X$ 上的向量丛就是向量空间。对一维丛，粘合函数 $f:(x,\theta)\mapsto\Aut V_x$ 就是连续映射 $f:S^1\to\C^*$。于是对每个 $n\in\Z$，用 $V=\C$、$f(\theta)=e^{in\theta}$ 作粘合构造，得到 $S^2$ 上的丛 $E_n$。注意由构造立刻得到 $\alpha[E_n]=-n$。回顾每个这样的映射 $f$ 都同伦于形如 $f(\theta)=e^{in\theta}$ 的映射，可见这一构造给出 $S^2$ 上所有一维丛。特别地，这说明 $\tilde K(S^2)\iso\Z$。下面从这一观点看 $S^2$ 上向量丛的几个例子。

\begin{example}
把 $S^2\iso\C\cup\infty$ 看成 Riemann 球面，考虑复切丛 $T_\C S^2$。把 $B^-$ 等同于单位圆盘 $\{z\mid |z|\le1\}$，把 $B^+$ 等同于 $\{w\mid |w|\le1\}$，其中 $w=z^{-1}$（处理复结构时，希望坐标转移是复解析的）。点 $z\in B^-$ 处的切向量可写成 $\xi=d/dt\,f(t)|_{t=0}$，其中 $f(0)=z$。换到 $w$ 坐标，相应切向量为
\[
\xi_w=\frac{d}{dt}f(t)^{-1}\Big|_{t=0}=-f(t)^{-2}\frac{d}{dt}f(t)\Big|_{t=0}=-z^{-2}\xi_z.
\]
于是可以写
\[
T_\C S^2\iso\bigl((B^-\times\C)\cup(B^+\times\C)\bigr)\big/\sim,
\]
其中在单位圆 $\{z\mid |z|=1\}=B^-\cap B^+$ 上有等同
\[
(z,\xi)\in B^-\times\C\ \sim\ (z,-z^{-2}\xi)\in B^+\times\C.
\]
把单位圆转半圈可以看出，$f(z)=-z^{-2}$ 同伦于映射 $f(x)=z^{-2}$。这个映射对应 $f(\theta)=e^{-2i\theta}$，故 $T_\C S^2\iso E_{-2}$。
\end{example}

\begin{example}
考虑映射 $f(\theta)=e^{in\theta}$。相应椭圆伪微分算子 $P_n$ 的象征 $\sigma(P_n)\in K(TS^1)$ 由 $TS^1$ 上的复形
\[
0\longrightarrow TS^1\times\C\xrightarrow{\ \sigma_n\ }TS^1\times\C\longrightarrow 0
\]
给出，其中同态由
\[
\sigma_n(\theta,\xi)=\begin{cases} e^{in\theta} & \xi>0,\\ 1 & \xi<0 \end{cases}
\]
给出。设 $p^+\in B^+$、$p^-\in B^-$ 是球面 $S^2$ 的「北极」与「南极」。可把 $B^--\{p^-\}$ 等同于 $\{(\theta,\xi)\in TS^1\mid\xi\le0\}$，把 $B^+-\{p^+\}$ 等同于 $\{(\theta,\xi)\in TS^1\mid\xi\ge0\}$。注意 $\sigma_n(\theta,\xi)$ 同伦于同态 $\sigma'_n(\theta,\xi)$，后者由
\[
\sigma'(\theta,\xi)=\begin{cases} f(\theta)/(1+\xi) & \xi>0,\\ 1 & \xi<0 \end{cases}
\]
给出。可以把 $\sigma'_n$ 延拓为连续映射 $\sigma'_n:E_0\to E_n$，只需令
\[
\sigma'_n(p^-)=\lim_{\xi\to-\infty}\sigma'_n(\theta,\xi)=1,\qquad
\sigma'_n(p^+)=\lim_{\xi\to+\infty}\sigma'_n(\theta,\xi)=0.
\]
注意这个映射在圆 $S^1\times\{0\}$ 上连续：按粘合构造，我们把点 $((\theta,0),z)\in B^-\times\C$ 与 $((\theta,0),e^{in\theta}z)\in B^+\times\C$ 等同。于是象征复形延拓为 $S^2$ 上的复形
\[
0\longrightarrow E_0\xrightarrow{\ \sigma'_n\ }E_n\longrightarrow 0,
\]
并且 $\sigma(P_n)=[E_0]-[E_n]\in\tilde K(S^2)$。由此得到 $\alpha(\sigma(P_f))=n$。
\end{example}

% ===== 原书 p.43 =====
\begin{proposition}
对任意 $x\in K(TS^1)\iso K(\R^2)\iso\tilde K(S^2)$，有 $\aind_{S^1}(x)=-\alpha_\pt(x)$。
\end{proposition}

\begin{proof}
注意 $K(TS^1)\iso\Z$，故 $K(TS^1)$ 中每个元素 $x$ 都作为某个椭圆伪微分算子 $P_n$ 的象征 $x\in\sigma(P_n)$ 出现。于是由上面的讨论
\[
\begin{aligned}
\aind_{S^1}(x)&=\Ind P_n=-n\\
&=-\alpha_\pt\bigl(\sigma(P_n)\bigr)=-\alpha_\pt(x).
\end{aligned}
\]
\end{proof}

\begin{example}
回顾 §1：Bott 类 $\lambda_1\in K(\C)\iso\tilde K(S^2)$ 由 $\C$ 上的外复形 $\Lambda(\C)$ 表示：
\[
0\longrightarrow\pi^*\Lambda^0(\C)\xrightarrow{\ \alpha\ }\pi^*\Lambda^1(\C)\longrightarrow 0,
\]
其中 $\pi^*\Lambda^0(\C)\iso\pi^*\Lambda^1(\C)\iso\C\times\C$，且 $\alpha:(z,v)\mapsto(z,zv)$。记 $D^-=\{z\mid |z|\le1\}$、$D^+=\{z\mid |z|\ge1\}$。通过包含 $\C\subset\C^+\iso S^2$，可把 $D^-$ 与 $B^-$ 等同，把 $D^+\cup\infty$ 与 $B^+$ 等同。因为 $E_0$ 是 $S^2$ 上的平凡丛，所以 $\pi^*\Lambda^0(\C)=E_0|_\C$。又可构造同构 $\beta:\pi^*\Lambda^1(\C)\to E_{-1}|_\C$，取
\[
\beta:(z,v)\mapsto\begin{cases}(z,v)\in B^-\times\C & z\in D^-,\\(z,z^{-1}v)\in B^+\times\C & z\in D^+.\end{cases}
\]
把圆周 $S^1=\R/2\pi\Z$ 与 $\C$ 中的单位圆等同，注意映射 $f(\theta)=e^{-i\theta}$ 等价于映射 $f(z)=z^{-1}$，故由粘合构造知 $\beta$ 连续。于是外复形 $\Lambda(\C)$ 可改写成 $\C$ 上的复形
\[
0\longrightarrow E_0|_\C\xrightarrow{\ \gamma\ }E_{-1}|_\C\longrightarrow 0,
\]
其中 $\gamma=\beta\circ\alpha$ 由
\[
\gamma:(z,v)\mapsto\begin{cases}(z,zv)\in B^-\times\C & z\in D^-,\\(z,v)\in B^+\times\C & z\in D^+.\end{cases}
\]
给出。最后令 $\gamma(\infty,v)=(\infty,v)$，可把它延拓为 $S^2$ 上的复形
\[
0\longrightarrow E_0\xrightarrow{\ \gamma\ }E_{-1}\longrightarrow 0,
\]
而 $S^2$ 紧，故 $\lambda_1=[E_0]-[E_{-1}]\in\tilde K(S^2)=K(\R^2)$。计算 $\alpha(\lambda_1)=\alpha[E_0]-\alpha[E_{-1}]=0-1=-1$，上面的命题给出 $\aind_{S^1}(\lambda_1)=-\alpha(\lambda_1)=1$。这正是 §3 中完成归一化公理证明所需的结果。
\end{example}

\transnote{符号约定要看清：由构造 $\alpha[E_n]=-n$，而 $\Ind P_n=-n$，两者一致；生成元 $b=[E_{-1}]-[E_0]$ 给出 $\alpha(b)=0-(-1)=1$。后文 Bott 周期性定理的 $\alpha(b)=1$ 与归一化公理都依赖这一符号。若把 $\alpha$ 改成取 $+\deg f$，整条链上的符号会同时反转。}

\subsection{Bott 周期性定理}

取 $b=-\lambda_1=[E_{-1}]-[E_0]$ 作为 $K(\R^2)$ 的生成元，于是 $\alpha(b)=1$。还注意到，由构造有交换图表
\[
\begin{tikzcd}[row sep=large, column sep=large]
K(\R^2\times X)\otimes K(Y) \arrow[r,"\BT"] \arrow[d,"\alpha_x\otimes1"'] & K(\R^2\times X\times Y) \arrow[d,"\alpha_{X\times Y}"] \\
K(X)\otimes K(Y) \arrow[r,"\BT"] & K(X\times Y)
\end{tikzcd}
\]

% ===== 原书 p.44 =====
它说明：对任意 $x\in K(\R^2\times X)$ 与 $y\in K(Y)$，有乘积公式
\[
\alpha(X)\cdot y=\alpha(x\cdot y).
\]
还注意到，映射 $\alpha:K(\R^2\times X)\to K(X)$ 可通过单点紧化 $X^+$ 扩充到局部紧空间 $X$，乘积公式仍成立。现在，Bott 周期性定理可由 K 理论的几条简单代数性质直接推出。

\begin{theorem}[Bott 周期性]
设 $X$ 是局部紧空间，则如下两个同态
\[
\begin{gathered}
\beta:K(X)\to K(\R^2\times X)\\
\alpha:K(\R^2\times X)\to K(X)
\end{gathered}
\]
互为逆。
\end{theorem}

\begin{proof}
先由乘积公式，对任意 $x\in K(X)$ 有
\[
\alpha\beta(X)=\alpha(b\cdot x)=\alpha(b)\cdot x=1\cdot x=x,
\]
故 $\alpha$ 是 $\beta$ 的左逆。

还要证明 $\alpha$ 也是 $\beta$ 的右逆，即对每个 $u\in K(\R^2\times X)$ 验证 $b\cdot\alpha(u)=\beta\alpha(u)=u$。由乘积公式有 $\alpha(u\cdot b)=\alpha(u)\cdot b$，其中 $u\cdot\lambda_1\in K(\R^2\times X\times\R^2)$，我们希望交换这两个乘积。为此考虑映射
\[
\tau:\R^2\times X\times\R^2\to\R^2\times X\times\R^2,
\]
它交换两个 $\R^2$ 的次序。取 $X$ 为一点，相应的变换 $\tau:\R^4\to\R^4$ 由矩阵
\[
\begin{pmatrix}0&I_2\\I_2&0\end{pmatrix},\qquad\text{其中 } I_2=\begin{pmatrix}1&0\\0&1\end{pmatrix}
\]
给出；其行列式为 $+1$，故它与单位元在 $\GL(4,\R)$ 的同一连通分支中。对一般的 $X$，同样可以构造 $\tau$ 与恒等映射之间的同伦，从而 K 理论上的诱导映射
\[
\tau^*:K(\R^2\times X\times\R^2)\to K(\R^2\times X\times\R^2)
\]
是恒等映射。设 $\tilde u\in K(X\times\R^2)$ 是 $u$ 在典范同构 $K(\R^2\times X)\to K(X\times\R^2)$ 下的对应元素，得到
\[
u\cdot b=\tau^*(u\cdot b)=b\cdot\tilde u,
\]
两边作用 $\alpha$，得
\[
\alpha(u)\cdot b=\alpha(u\cdot b)=\alpha(b\cdot\tilde u)=\tilde u\in K(X\times\R^n).
\]
再经由同构 $K(X\times\R^2)\to K(\R^2\times X)$ 变回去，得到
\[
u=b\cdot\alpha(u)=\beta\alpha(u)\in K(\R^2\times X),
\]
证明完成。
\end{proof}

\transnote{关键是 $\det\tau=+1$：两块 $\R^2$ 互换的矩阵行列式为 $\det(-I_2)=+1$，故 $\tau$ 落在 $\GL(4,\R)$ 的单位分支，$\tau^*$ 才是恒等映射，两个乘积才能交换。维数 $2$ 在此处省不掉：若把两个 $\R^2$ 换成两个 $\R^1$，交换矩阵行列式为 $-1$，$\tau^*$ 不再平凡，$\alpha$ 只能是右逆的反号版本。原文先取 $X$ 为一点算行列式，再靠对 $X$ 的自然性回到一般情形。}

% ===== 原书 p.45 =====
\subsection{Thom 同构定理}

这一 Bott 周期性证明因构造是「模块化」的，容易推广。事实上，我们只需要知道逆同态 $\alpha_X:K(\R^2\times X)\to K(X)$ 满足下面两条初等性质：

\textbf{乘积公式}\quad 对一切 $x\in K(\R^2\times X)$ 与 $y\in K(Y)$，$\alpha_{X\times Y}(x\cdot y)=\alpha_X(x)\cdot y$。

\textbf{归一化公理}\quad $\alpha_\pt(b)=1$，其中 $b=\beta_\pt(1)\in K(\R^2)$。

因此，若换一种构造给出逆同态 $\alpha_X:K(\R^2\times X)\to K(X)$，只要它仍满足上述两条性质，我们的证明依然成立。另一种这样的构造见 [7]。

现在把 Bott 周期性的证明推广到等变情形。设给定紧 Lie 群 $G$。取有限维复 $G$-模 $V$，可构造元素 $\lambda_V\in K_G(V)$，它由外复形 $\Lambda(V)$ 给出。把讨论中的 $\R^2$ 换成 $V$，得到同态
\[
\varphi:K_G(X)\xrightarrow{\ \lambda_V^*\BT\cdot\ }K_G(V\times X),
\]
它由 Thom 类 $\lambda_V$ 的对偶丛 $\lambda_V^*$ 作外乘给出。要证这个映射是同构，须构造逆同态 $\alpha_X:K_G(V\times X)\to K_G(X)$，使它满足如下两条类似公理：

\textbf{乘积公式}\quad 对一切 $x\in K_G(V\times X)$ 与 $y\in K_G(Y)$，$\alpha_{X\times Y}(x\cdot y)=\alpha_X(x)\cdot y$。

\textbf{归一化公理}\quad $\alpha_\pt(\lambda_V^*)=1\in R(G)$，其中 $\lambda_V^*=\varphi_\pt(1)\in K_G(V)$。

这样的构造同样可以用一族椭圆算子的指标做到（见 [4, §4]），于是得到

\begin{theorem}[等变周期性]
若 $X$ 是紧 $G$-空间，$V$ 是有限维复 $G$-模，则由 $\lambda_V^*\in K_G(V)$ 外乘给出的同态
\[
\varphi:K_G(X)\to K_G(V\times X)
\]
是同构。
\end{theorem}

本定理有一个特例：取 $G=\U(n)$ 为酉群，$V=\C^n$ 为平凡表示。此时 $X$ 上每个复向量丛都可写成 $E=Y\times_{\U(n)}\C^n$，其中 $Y$ 是自由 $\U(n)$-空间，$Y/\U(n)=X$。回顾 §4：若 $X$ 是自由 $G$-空间，则 $K_G(X)\iso K(X/G)$。于是
\[
K_{\U(n)}(Y)\iso K(X),\qquad K_{\U(n)}(Y\times_{\U(n)}\C^n)\iso K(E).
\]
注意 Thom 类 $\lambda_E\in K(E)$ 对应元素 $\varphi(\lambda_V)\in K_G(V\times X)$，由等变周期性定理立刻得到

% ===== 原书 p.46（上半页，仍属附录 A） =====
\begin{corollary}[Thom 同构定理]
若 $E$ 是紧空间 $X$ 上的复向量丛，则用 Thom 类 $\lambda_E\in K(E)$ 作乘法诱导出同构 $K(X)\to K(E)$。
\end{corollary}

还可以引入第二个群 $H$ 的作用，讨论不受影响。若 $X$ 是 $H$-空间，$E$ 是 $X$ 上的 $H$-丛，则可把 $E$ 写成 $E=Y\times_{\U(n)}\C^n$，其中 $Y$ 是 $H\times\U(n)$-空间，$Y/\U(n)=X$。于是
\[
K_{H\times\U(n)}(Y)\iso K_H(X),\qquad K_{H\times\U(n)}(Y\times_{\U(n)}\C^n)\iso K_H(E),
\]
从而得到等变推广。

\begin{corollary}[等变 Thom 同构定理]
若 $E$ 是紧 $H$-空间 $X$ 上的复 $H$-向量丛，则用 Thom 类 $\lambda_E\in K_H(E)$ 作乘法诱导出同构 $K_H(X)\to K_H(E)$。
\end{corollary}

% ===== 原书 p.46（下半页，附录 B 起） =====
\section{Atiyah--Hirzebruch 谱序列}

这里给出计算 K 理论的 Atiyah--Hirzebruch 谱序列，叙述依 [8, §2]。我们不假定读者熟悉由 Cartan--Eilenberg 系统构造谱序列（见 [14]），而是从对与三元组的长正合列出发，推出所需的全部同调代数。

设 $X$ 是有限（连通）CW 复形，$X^p$ 记它的 $p$-骨架。定义
\[
K^*_p(X)=\Ker\{K^*(X)\to K^*(X^{p-1})\}.
\]
这就是 K 理论的 Atiyah--Hirzebruch 滤过，且
\[
K^*_0(X)=K^*(X)\supset K^*_1(X)=\tilde K^*(X)\supset\cdots\supset K^*_{\max}(X)=0,
\]
其中 $X^{-1}=\varnothing$，$X^{\max-1}=X$。\footnote{尽管从定义看不出，Atiyah--Hirzebruch 滤过其实是同伦型不变量，与 $X$ 上 CW 结构的选取无关。此外这一滤过是乘性的，即 $K^*_p(X)\cdot K^*_q(X)\subset K^*_{p+q}(X)$ 对一切 $p,q$ 成立。更多信息见 [8, §2]。} Atiyah--Hirzebruch 谱序列计算相伴分次分量
\[
\operatorname{Gr}_p K^*(X)=K^*_p(X)\big/K^*_{p+1}(X).
\]
若知道这些分次分量，就可以用一串群扩张去还原 $K^*(X)$。

现在考虑如下正合列的交换图表：
\[
\begin{tikzcd}[column sep=small, row sep=large]
K^{*-1}(X^p) \arrow[r] \arrow[d] & K^*(X,X^p) \arrow[r] \arrow[d] & K^*(X) \arrow[r] \arrow[d,equal] & K^*(X^p) \arrow[d] \\
K^{*-1}(X^{p-1}) \arrow[r] \arrow[d] & K^*(X,X^{p-1}) \arrow[r] \arrow[d] & K^*(X) \arrow[r] & K^*(X^{p-1}) \\
K^*(X^p,X^{p-1}) \arrow[r,equal] \arrow[d] & K^*(X^p,X^{p-1}) \arrow[d] & & \\
K^*(X^p) \arrow[r] & K^{*+1}(X,X^p) & &
\end{tikzcd}
\]
其中上两行是 $(X,X^p)$ 与 $(X,X^{p-1})$ 的长正合列，左两列是 $(X^p,X^{p-1})$ 与三元组 $(X,X^p,X^{p-1})$ 的长正合列。

\transnote{图表中两处双线含义不同：$K^*(X)$ 处的双线只是把上两行的同一列对齐，标出两行长正合列在此共用该群；$K^*(X^p,X^{p-1})$ 处的双线则表示同一个相对群。两者都不是同构箭头，读图时不必找对应的映射。}

% ===== 原书 p.47 =====
由 $K^*_p(X)$ 的定义与上图表上两行的正合性，在 $K^*(X)$ 中有
\[
\begin{aligned}
K^*_p(X)&=\operatorname{Im} K^*(X,X^{p-1})\iso K^*(X,X^{p-1})\big/\operatorname{Im} K^{*-1}(X^{p-1}),\\
K^*_{p+1}(X)&=\operatorname{Im} K^*(X,X^p)\iso K^*(X,X^p)\big/\operatorname{Im} K^{*-1}(X^p),
\end{aligned}
\]
于是相伴分次分量为
\[
\begin{aligned}
\operatorname{Gr}_p K^*(X)&\iso K^*(X,X^{p-1})\big/\bigl(\operatorname{Im} K^*(X,X^p)+\operatorname{Im} K^{*-1}(X^{p-1})\bigr)\\
&\iso\bigl(K^*(X,X^{p-1})\big/\operatorname{Im} K^*(X,X^p)\bigr)\big/\operatorname{Im} K^{*-1}(X^{p-1})\\
&\iso\Ker\{K^*(X^p,X^{p-1})\to K^{*+1}(X,X^p)\}\big/\operatorname{Im} K^{*-1}(X^{p-1}).
\end{aligned}
\]
为计算这个分次分量，构造一个谱序列，其第一页为
\[
E^p_1:=K^*(X^p,X^{p-1}),
\]
第 $k$ 页由
\[
E^p_k=\frac{\Ker\{K^*(X^p,X^{p-1})\xrightarrow{\ \beta\ }K^{*+1}(X^{p-1+k},X^p)\}}
{\operatorname{Im}\{K^{*-1}(X^{p-1},X^{p-k})\xrightarrow{\ \beta'\ }K^*(X^p,X^{p-1})\}}
\]
给出。注意对 $E_1$，从
\[
E^p_1=\Ker\{K^*(X^p,X^{p-1})\to K^{*+1}(X^p,X^p)\}\big/\operatorname{Im} K^{*-1}(X^{p-1},X^{p-1})
\]
出发，这与 $E^p_1$ 的定义一致。另一方面，$k$ 足够大时 $X^{p-1+k}=X$、$X^{p-k}=\varnothing$，于是序列稳定下来，给出所需的分次分量 $E^p_\infty=\operatorname{Gr}_p K^*(X)$。考察下面的交换图表（取自前一图表的左下角），
\[
\begin{tikzcd}[row sep=large, column sep=large]
K^{*-1}(X^{p-1},X^{p-k}) \arrow[r] \arrow[d] & K^*(X^{p-1+k},X^{p-1}) \arrow[d] \\
K^*(X^p,X^{p-1}) \arrow[r,equal] & K^*(X^p,X^{p-1}) \arrow[d] \\
& K^{*+1}(X^{p-1+k},X^p)
\end{tikzcd}
\]
可见在 $K^*(X^p,X^{p-1})$ 中
\[
\begin{aligned}
\operatorname{Im} K^{*-1}(X^{p-1},X^{p-k})&\subset\operatorname{Im} K^{*-1}(X^{p-1+k},X^{p-1})\\
&=\Ker\{K^*(X^p,X^{p-1})\to K^{*+1}(X^{p-1+k},X^p)\},
\end{aligned}
\]
从而 $E^p_k$ 确实良定义。

\transnote{稳定要两个条件同时成立：$X^{p-1+k}=X$ 需 $p-1+k\ge\max$，$X^{p-k}=\varnothing$ 需 $p-k\le-1$。取 $k\ge\max-p+1$ 即可，故每个 $p$ 上 $E^p_k$ 终将稳定，$\max$ 是最高非空骨架的序号。}

% ===== 原书 p.48 =====
为构造谱序列的微分，考虑如下交换图表：
\[
\begin{tikzcd}[row sep=large, column sep=small]
& K^{*-1}(X^{p-1},X^{p-k}) \arrow[r] \arrow[dr,"\beta'"'] & K^{*-1}(X^{p-1},X^{p-k-1}) \arrow[d,"\alpha'"] & \\
& & K^*(X^p,X^{p-1}) \arrow[dll,dashed,"d_k"'] \arrow[d,"\alpha"] \arrow[dr,"\beta"] & \\
K^*(X^{p-1+k},X^p) \arrow[r,"\epsilon"] & K^{*+1}(X^{p+k},X^{p-1+k}) \arrow[r,"\delta"] & K^{*+1}(X^{p+k},X^p) \arrow[r,"\gamma"] & K^{*+1}(X^{p-1+k},X^p)
\end{tikzcd}
\]
其中上行来自三元组 $(X^{p-1},X^{p-k},X^{p-k-1})$ 的长正合列，下行来自三元组 $(X^{p+k},X^{p-1+k},X^p)$ 的长正合列。目标是：用群 $E^p_k=\Ker\beta/\operatorname{Im}\beta'$ 计算群 $E^{p+k}_{k+1}=\Ker\alpha/\operatorname{Im}\alpha'$。

设 $x\in K^*(X^p,X^{p-1})$ 是类 $[x]\in E^p_k$ 的代表元。由 $E^p_k$ 的定义 $x\in\Ker\beta$，于是由交换图表 $\alpha x\in\Ker\gamma=\operatorname{Im}\delta$。因此可写 $\alpha x=\delta y$，其中 $y\in K^{*+1}(X^{p+k},X^{p-1+k})$，且 $y$ 在 $\Ker\delta=\operatorname{Im}\epsilon$ 的意义下确定。这样就得到良定义的类
\[
[y]\in K^{*+1}(X^{p+k},X^{p-1+k})\Big/\operatorname{Im}\Bigl\{K^*(X^{p-1+k},X^p)\xrightarrow{\ \epsilon\ }K^{*+1}(X^{p+K},X^{p-1+k})\Bigr\},
\]
并把谱序列的微分定义为
\[
d_k:E^p_k\to E^{p+k}_k,\qquad d_k[x]:=[y].
\]
为完备起见，还应验证
\[
y\in\Ker\{K^{*+1}(X^{p+k},X^{p-1+k})\to K^{*+2}(X^{p+2k-1},X^{p+k})\}
\]
（用类似的图表追踪），验证类 $d_k[x]\in E^{p+k,*+1}_k$ 与 $x$ 的选取无关（因为 $\operatorname{Im}\beta'\subset\operatorname{Im}\alpha'\subset\Ker\alpha$），并验证 $d_k$-上同调确实给出
\[
\Ker d_k/\operatorname{Im} d_k=\Ker\alpha/\operatorname{Im}\alpha'=E^p_{k+1},
\]
即谱序列的下一页。

为把上同调分次并入上述讨论，在谱序列上使用双分次：项 $E^{p,q}_k$ 带的上同调次数为 $p+q$。微分于是给出
\[
d_k:E^{p,q}_k\to E^{p+k,q-k+1}_k
\]
对一切 $p,q$ 成立。这个谱序列的 $E_1$ 页有分量 $E^{p,q}_1=K^{p+q}(X^p,X^{p-1})$。然而，若相对 $(p-1)$-骨架来作，则 $X^p$ 成为一束 $p$-球面（bouquet），每个 $p$-胞腔对应一个。于是得到
\[
E^{p,q}_1=K^{p+q}(X^p,X^{p-1})\iso K^{p+q}(S^p)^{\oplus\beta_p}\iso K^q(\pt)^{\oplus\beta_p}=C^p\bigl(X;K^q(\pt)\bigr)
\]
（其中 $\beta_p$ 是 $p$-胞腔的个数），这就是 $X$ 的、以一点的 K 理论为系数的胞腔上同调的链群。此外，第一个微分
\[
d_1:C^p\bigl(X;K^q(\pt)\bigr)\to C^{p+1}\bigl(X;K^q(\pt)\bigr)
\]

% ===== 原书 p.49 =====
正是胞腔链复形中的微分（见 [2, §III.7]），于是
\[
E^{p,q}_2=H^p\bigl(X;K^q(\pt)\bigr)
\]
就是 Atiyah--Hirzebruch 谱序列的起点。

上述全部讨论不仅对 K 理论成立，对任何广义上同调理论也成立。在复 K 理论这一具体情形，事情还会进一步简化。得到
\[
E^{p,q}_2=\begin{cases}H^p(X;\Z) & q \text{ 偶},\\ 0 & q \text{ 奇},\end{cases}
\]
于是整个谱序列只需用 $q=0$ 那一行计算。$k$ 为偶数的微分 $d_k$ 全为零，因为它们交换奇偶 $q$。$k$ 为奇数时，微分为上同调运算：
\[
d_k:H^p(X;\Z)\to H^{p+k}(X;\Z).
\]
因此，计算复 K 理论从整系数上同调开始。然后依次对 $k$ 为奇数的运算 $d_k$ 取上同调。最后，把偶与奇上同调群剩余部分用一串（常常非平凡的）群扩张组合起来，还原 $K^0(X)$ 与 $K^1(X)$。这个论证的一个推论是：K 理论从不超过整系数上同调。

\transnote{「$q$ 偶时为 $H^p(X;\Z)$、$q$ 奇时为零」只在复 K 理论成立，依据是 $K^q(\pt)$ 随 $q$ 的奇偶取 $\Z$ 或 $0$。实 K 理论的 $KO^q(\pt)$ 以 $8$ 为周期，$E_2$ 页并非只有一行，$d_k$ 也不再只作用在一行上。}

\begin{example}
考虑实射影空间 $X=\R P^n$。它的整系数上同调为
\[
H^*(\R P^n;\Z)\iso\begin{cases}
\Z & \text{次数 }0,\\
\Z/2\Z & 0 \text{ 与 } n \text{ 之间的偶数次数},\\
\Z & n \text{ 为奇数时的次数 } n,\\
0 & \text{其余情形}.
\end{cases}
\]
次数 $0$ 的分量 $\Z$ 必须留存到 K 理论，因为它对应一点的上同调或 K 理论（这个 $\Z$ 确实不出现在约化上同调与约化 K 理论中），故所有微分 $d_k$ 在它上面为零。此外不存在非零同态 $\Z/2\Z\to\Z$，故最高次分量 $\Z$（$n$ 为奇数时）总落在微分 $d_k$ 的余核中。$\Z/2\Z$ 分量只出现在偶数次数，故所有奇数微分 $d_k$ 为零；偶数微分又总为零。于是 Atiyah--Hirzebruch 谱序列在 $E_2$ 页退化。由此 $K^*(\R P^n)$ 在次数 $0$ 有一个 $\Z$ 分量，$n$ 为奇数时在次数 $1$ 有一个 $\Z$ 分量，次数 $0$ 还有由一串 $\Z/2\Z$-扩张得到的挠分量。但这些信息没有完全确定 $\R P^n$ 的 K 理论：每个 $\Z/2\Z$-扩张都有两种可能形式：
\[
0\to\Z/2\Z\to\Z/4\Z\to\Z/2\Z\to0,
\]
或
\[
0\to\Z/2\Z\to\Z/2\Z\oplus\Z/2\Z\to\Z/2\Z\to0.
\]
事实上，$K(\R P^n)$ 用前者（见 [6, Proposition 2.7.7]）：
\[
K^0(\R P^n)\iso\Z\oplus\Z/2^k\Z,\qquad
K^1(\R P^n)\iso\begin{cases}\Z & n \text{ 奇},\\ 0 & n \text{ 偶},\end{cases}
\]
其中 $k=2^{\lfloor n/2\rfloor}$。一串 $\Z/2\Z$ 分次分量经由群扩张逐次组合，得到 $\Z/2\Z$、$\Z/4\Z$、$\Z/8\Z$ 等等。
\end{example}

\transnote{$d_k$ 的下标 $k$ 与 $K^0(\R P^n)\iso\Z\oplus\Z/2^k\Z$ 中的 $k$ 是两回事：前者是谱序列页码，后者是循环群的指数，此处 $k=2^{\lfloor n/2\rfloor}$。同一符号在相邻两处表示不同对象，是本节最易读错的地方。}

% ===== 原书 p.50 =====
\begin{thebibliography}{99}
\bibitem{1} J. F. Adams. \emph{Lectures on Lie groups}. W. A. Benjamin, Inc., New York-Amsterdam, 1969.
\bibitem{2} J. F. Adams. \emph{Stable homotopy and generalised homology}. Chicago Lectures in Mathematics. University of Chicago Press, Chicago, IL, 1995.
\bibitem{3} M. F. Atiyah. Algebraic topology and elliptic operators. \emph{Comm. Pure Appl. Math.}, 20:237--249, 1967.
\bibitem{4} M. F. Atiyah. Bott periodicity and the index of elliptic operators. \emph{Quart. J. Math. Oxford Ser. (2)}, 19:113--140, 1968.
\bibitem{5} M. F. Atiyah. Algebraic topology and operators in Hilbert space. In \emph{Lectures in Modern Analysis and Applications. I}, pages 101--121. Springer, Berlin, 1969.
\bibitem{6} M. F. Atiyah. \emph{K-theory}. Advanced Book Classics. Addison-Wesley Publishing Company Advanced Book Program, 1989.
\bibitem{7} M. F. Atiyah and R. Bott. On the periodicity theorem for complex vector bundles. \emph{Acta Math.}, 112:229--247, 1964.
\bibitem{8} M. F. Atiyah and F. Hirzebruch. Vector bundles and homogeneous spaces. In \emph{Proc. Sympos. Pure Math., Vol. III}, pages 7--38. American Mathematical Society, Providence, R.I., 1961.
\bibitem{9} M. F. Atiyah and I. M. Singer. The index of elliptic operators on compact manifolds. \emph{Bull. Amer. Math. Soc.}, 69:422--433, 1963.
\bibitem{10} M. F. Atiyah and I. M. Singer. The index of elliptic operators. I. \emph{Ann. of Math. (2)}, 87:484--530, 1968.
\bibitem{11} M. F. Atiyah and I. M. Singer. The index of elliptic operators. III. \emph{Ann. of Math. (2)}, 87:546--604, 1968.
\bibitem{12} R. Bott. The index theorem for homogeneous differential operators. In \emph{Differential and Combinatorial Topology (A Symposium in Honor of Marston Morse)}, pages 167--186. Princeton University Press, Princeton, N.J., 1965.
\bibitem{13} R. Bott and L. W. Tu. \emph{Differential forms in algebraic topology}, volume 82 of Graduate Texts in Mathematics. Springer-Verlag, New York, 1982.
\bibitem{14} H. Cartan and S. Eilenberg. \emph{Homological Algebra}. Princeton Landmarks in Mathematics. Princeton University Press, Princeton, N.J., 1999.
\bibitem{15} F. Hirzebruch. Lectures on K-theory. In J. F. Adams, editor, \emph{Algebraic Topology: A Student's Guide}, volume 4 of London Mathematical Society Lecture Notes Series, pages 223--238. Cambridge University Press, London, 1972. Mimeographed lecture notes of the AMS Summer Topology Institute, Seattle, 1963.
\bibitem{16} J. W. Milnor and J. D. Stasheff. \emph{Characteristic classes}. Princeton University Press, Princeton, N.J., 1974.
\bibitem{17} G. Segal. Equivariant K-theory. \emph{Inst. Hautes \'Etudes Sci. Publ. Math.}, 34:129--151, 1968.
\end{thebibliography}


\end{document}
